\documentclass[11pt,letterpaper]{report}

% Edición en español, sincronizada con la revisión de manuscrito v0.5.0 (inglés: ../main.tex).
\usepackage[utf8]{inputenc}
\DeclareUnicodeCharacter{2212}{\ensuremath{-}}% U+2212 minus sign used in the abstracts
\usepackage[T1]{fontenc}
\usepackage{mathptmx}
\usepackage{microtype}
\usepackage[letterpaper,left=3cm,right=2.5cm,top=2.7cm,bottom=2.5cm]{geometry}
\usepackage{amsmath,amssymb}
\usepackage{graphicx}
\usepackage{booktabs,tabularx,array,multirow,makecell}
\usepackage{enumitem}
\usepackage{setspace}
\usepackage[font=small,labelfont=bf,labelsep=period,position=top]{caption}
\usepackage{float}
\usepackage{tikz}
\usetikzlibrary{arrows.meta,positioning,calc}
\usepackage{pgfplots}
\pgfplotsset{compat=1.16}
\usepackage[numbers,sort&compress]{natbib}
\usepackage{titlesec}
\usepackage{fancyhdr}
\usepackage{xurl}
\usepackage{hyperref}

\hypersetup{
  colorlinks=true,
  linkcolor=black,
  citecolor=black,
  urlcolor=blue,
  pdflang={es-CO},
  pdftitle={Seguimiento de precios y fragilidad de liquidez en mercados cripto híbridos CEX/DEX},
  pdfauthor={Edward Calderón},
  pdfsubject={Tesis de maestría con contraste confirmatorio sintético de H1, H2 y H3, 24 de noviembre de 2026, manuscrito v0.5.0}
}

\newcommand{\languageCode}{es}
\newcommand{\versionTesis}{v0.1.7}
\newcommand{\revisionManuscrito}{v0.5.0}
\newcommand{\datadir}{../figures/data}
\newcommand{\RunID}{cex-dex-confirm-5b06486d695e}

% Nombres en español (no hay patrones de separación silábica en español en esta instalación).
\renewcommand{\contentsname}{Contenido}
\renewcommand{\listtablename}{Lista de tablas}
\renewcommand{\listfigurename}{Lista de figuras}
\renewcommand{\tablename}{Tabla}
\renewcommand{\figurename}{Figura}
\renewcommand{\bibname}{Bibliografía}
\renewcommand{\appendixname}{Anexo}
\renewcommand{\chaptername}{Capítulo}
\renewcommand{\refname}{Bibliografía}
\hyphenpenalty=10000
\exhyphenpenalty=10000
\tolerance=1600
\emergencystretch=3em

% Notación común a ambas ediciones.
\newcommand{\pD}{p^{\mathrm{D}}}
\newcommand{\pC}{p^{\mathrm{C}}}
\newcommand{\E}{\mathbb{E}}
\newcommand{\sgn}{\operatorname{sgn}}
\newcommand{\code}[1]{\texttt{#1}}

\setlength{\parskip}{0.35em}
\setlength{\parindent}{0pt}
\setlength{\headheight}{14pt}
\renewcommand{\arraystretch}{1.12}
\setlist{leftmargin=*,itemsep=0.15em,topsep=0.3em}
\setcounter{tocdepth}{2}
\numberwithin{equation}{chapter}
\captionsetup{skip=6pt}
\titleformat{\chapter}[hang]{\Large\bfseries}{\thechapter.}{0.6em}{}
\titlespacing*{\chapter}{0pt}{-8pt}{16pt}
\titleformat{\section}{\large\bfseries}{\thesection.}{0.5em}{}
\titleformat{\subsection}{\normalsize\bfseries}{\thesubsection.}{0.5em}{}
\renewcommand{\bibsection}{\chapter{\bibname}}
\setlength{\bibsep}{3pt}

\pagestyle{fancy}
\fancyhf{}
\fancyhead[L]{\small CEX/DEX: seguimiento de precios y liquidez}
\fancyhead[R]{\small Manuscrito \revisionManuscrito{} · 24 nov. 2026}
\fancyfoot[C]{\small\thepage}
\renewcommand{\headrulewidth}{0.3pt}
\fancypagestyle{plain}{%
  \fancyhf{}
  \fancyhead[L]{\small CEX/DEX: seguimiento de precios y liquidez}
  \fancyhead[R]{\small Manuscrito \revisionManuscrito{} · 24 nov. 2026}
  \fancyfoot[C]{\small\thepage}
  \renewcommand{\headrulewidth}{0.3pt}
}

\begin{document}

% ---------------------------------------------------------------- Portada
\begin{titlepage}
  \thispagestyle{empty}
  \centering
  \singlespacing
  {\large\bfseries Universidad Nacional de Colombia\\[2pt]
  Facultad de Minas\\[2pt]
  Sede Medellín\par}
  \vspace{3.2cm}
  {\LARGE\bfseries Seguimiento de precios y fragilidad de liquidez\\
  en mercados cripto híbridos CEX/DEX\par}
  \vspace{0.9cm}
  {\Large Un marco experimental basado en agentes\\
  y un modelo dinámico reducido\par}
  \vspace{2.6cm}
  {\Large Edward Calderón Aristizábal\par}
  \vspace{1.6cm}
  {\large Maestría en Ingeniería -- Ingeniería de Sistemas e Informática\\[2pt]
  Modalidad de profundización\par}
  \vspace{0.8cm}
  {\large Director: Prof. Juan David Ospina Arango\par}
  \vfill
  {\large Medellín, Colombia\\[2pt]
  24 de noviembre de 2026\par}
  \vspace{0.3cm}
  {\small Manuscrito de trabajo para revisión de la dirección --- revisión \revisionManuscrito{} (edición \versionTesis{})\par}
\end{titlepage}

\pagenumbering{roman}
\setcounter{page}{1}

% ---------------------------------------------------------------- Resumen / Abstract
\begingroup
\setcounter{secnumdepth}{-1}
\addtocontents{toc}{\protect\setcounter{tocdepth}{0}}
\titleformat{\section}{\Large\bfseries}{}{0pt}{}
\singlespacing
\section{Resumen}

Los mercados de criptoactivos combinan plataformas centralizadas (CEX) con creadores automáticos de mercado (AMM), cuyo precio depende de las reservas y de las operaciones ejecutadas. El arbitraje conecta ambos precios, pero su capacidad de corrección está limitada por costos, inventario y liquidez. Este trabajo estudia cómo un AMM de producto constante sigue un precio CEX exógeno después de una perturbación, en un micromundo con agentes de ruido, tendencia, arbitraje y provisión de liquidez, y evalúa un mapa dinámico discreto de dos variables como descripción reducida. Se contrastan tres hipótesis con un protocolo confirmatorio fijado de antemano: semillas y celdas de confirmación disjuntas de las de exploración, inferencia pareada por semilla y corrección de Holm.

Con 300 semillas de confirmación y 13500 ejecuciones, el arbitraje activo reduce el área de exceso de error de 2.36 a 0.10 (H1; diferencia media de −2.26; p de Holm 0.00015), y la capacidad gradúa la recuperación: 56.8 por ciento con capacidad baja y 100 por ciento con capacidad alta. El retiro adaptativo de liquidez aumenta en promedio la respuesta máxima en 0.011 (H2), pero solo con la menor profundidad; con profundidades mayores el efecto es nulo. El mapa reducido supera a la persistencia (H3), aunque no a un autorregresivo escalar, y su clasificación de regímenes fuera de muestra es débil. Un análisis suplementario de sensibilidad con 40\,500 trayectorias encontró que el 31.2 por ciento de las trayectorias adaptativas que se recuperaron salió después de la banda; la amplificación con profundidad baja desapareció al escalar las órdenes con las reservas y bajó de 0.0346 a 0.0147 unidades logarítmicas al desactivar el pánico. Estos resultados condicionados sustentan requisitos candidatos de monitoreo y pruebas de estrés, no prescripciones para mercados activos. Todos los hallazgos son sintéticos, carecen de validación externa y no se alcanzó la precisión objetivo de H1.

\textbf{Palabras clave:} simulación basada en agentes; arbitraje; liquidez; seguimiento de precios; mercados descentralizados; sistemas dinámicos discretos

\clearpage
\section{Abstract}

Cryptoasset markets combine centralized exchanges (CEX) with automated market makers (AMM), whose price depends on reserves and on the trades executed against them. Arbitrage links the two prices, but its capacity to correct them is limited by costs, inventory and liquidity. This study examines how a constant-product AMM tracks an exogenous CEX price after a shock, in a micro-world with noise, trend, arbitrage and liquidity-provision agents, and evaluates a two-variable discrete dynamical map as a reduced description. Three hypotheses are tested under a confirmatory protocol fixed in advance: confirmation seeds and cells disjoint from the exploration ones, inference paired by seed, and Holm correction.

With 300 confirmation seeds and 13500 runs, active arbitrage lowers the excess-error area from 2.36 to 0.10 (H1; mean difference of −2.26; Holm p of 0.00015), and capacity grades recovery: 56.8 percent with low capacity and 100 percent with high capacity. Adaptive liquidity withdrawal raises the maximum response by 0.011 on average (H2), but only at the shallowest depth; at greater depths the effect is null. The reduced map beats persistence (H3), though not a scalar autoregression, and its out-of-sample regime classification is weak. The results are synthetic, use no market data and have no external validation, and the target precision for H1 was not reached with 300 seeds.

\textbf{Keywords:} agent-based simulation; arbitrage; liquidity; price tracking; decentralized exchanges; discrete dynamical systems

\addtocontents{toc}{\protect\setcounter{tocdepth}{2}}
\endgroup
\setcounter{secnumdepth}{2}
\titleformat{\section}{\large\bfseries}{\thesection.}{0.5em}{}

\clearpage
\tableofcontents
\clearpage
\listoftables
\clearpage
\listoffigures
\clearpage
\chapter*{Notación}
\addcontentsline{toc}{chapter}{Notación}

Los símbolos siguientes se usan en todo el documento. El superíndice $\mathrm{D}$ designa al intercambio descentralizado (AMM) y $\mathrm{C}$ a la referencia centralizada.

\begin{center}
\small
\begin{tabularx}{\textwidth}{@{}l X@{}}
\toprule
Símbolo & Significado \\
\midrule
$x_t,\ y_t$ & Reservas de los activos $X$ e $Y$ del AMM en el tick $t$ (precio cotizado en unidades de $Y$ por unidad de $X$). \\
$k_t = x_t y_t$ & Producto de reservas; constante en un intercambio sin comisión. \\
$\pD_t = y_t/x_t$ & Precio marginal del AMM. \\
$\pC_t$ & Precio de referencia centralizado (exógeno). \\
$q$ & Cantidad de $X$ vendida al AMM en un intercambio. \\
$f$ & Comisión proporcional retenida por la piscina. \\
$I(q,x)$ & Impacto logarítmico de una venta: $2\log(1+q/x)$ sin comisión. \\
$e_t = \log(\pD_t/\pC_t)$ & Error logarítmico de seguimiento. \\
$\gamma$ & Semiancho logarítmico de la banda operativa de no arbitraje (valor ejecutado: $0.003$). \\
$d_t,\ D_t = |d_t|$ & Exceso de error fuera de la banda con signo, y su valor absoluto. \\
$c$ & Capacidad de arbitraje: fracción máxima de la reserva negociable por tick. \\
$q^{\ast}$ & Tamaño de venta que lleva el precio al borde de la banda. \\
$z_t$ & Presión neta de compra (ruido y tendencia; sin arbitraje) normalizada por la reserva inicial de $Y$. \\
$u_t = -\Delta\log \pC_t$ & Entrada exógena del mapa reducido. \\
$\alpha,\ \beta,\ \chi,\ \rho$ & Coeficientes del mapa de dos estados; $a = 1-\alpha$ y $c_z = -\chi$ en las tablas de ajuste. \\
$A$ & Matriz local del mapa autónomo; $\operatorname{tr}A$ y $\det A$ su traza y determinante. \\
$H$ & Horizonte de la ventana de respuesta (20 ticks). \\
$A_{\mathrm{ex}} = \sum_{t=1}^{H} D_t$ & Área de exceso de error (métrica primaria de H1). \\
$M = \max_{1\le t\le H}\bigl|\log(\pD_t/\pD_0)\bigr|$ & Respuesta máxima del precio del AMM (métrica primaria de H2). \\
$\Delta_1,\Delta_2,\Delta_3$ & Contrastes pareados de H1, H2 y H3. \\
$\mathrm{MAE}$ & Error absoluto medio de predicción a un paso (métrica primaria de H3). \\
E0--E3 & Escenarios experimentales (tabla~\ref{tab:escenarios}). \\
OE1--OE4 & Objetivos específicos. \\
\bottomrule
\end{tabularx}
\end{center}

\onehalfspacing

% ---------------------------------------------------------------- Cuerpo
\clearpage
\pagenumbering{arabic}
\setcounter{page}{1}
\chapter{Introducción}
\label{cap:introduccion}

Los criptoactivos se negocian en intercambios centralizados (CEX) y descentralizados (DEX), incluidos los creadores automáticos de mercado (AMM), cuyos precios cotizados dependen de las reservas de la piscina. El arbitraje puede acercar el precio de un AMM a una referencia externa, pero la magnitud y velocidad del ajuste dependen de la regla de intercambio, los costos, las reservas disponibles, el capital y el orden de las transacciones. Estos mecanismos se estudian en la literatura sobre pérdida frente al rebalanceo (LVR), arbitraje y valor extraíble por mineros o máximo (MEV) \cite{milionis2022lvr,daian2020flash,zhou2021hft}. Esta tesis aísla un entorno mucho más acotado: una piscina de producto constante sin comisiones, una referencia CEX exógena y una regla de arbitraje limitado.

La primera contribución es analítica. Un intercambio de $c$ veces la reserva de entrada cambia el precio logarítmico en $2\log(1+c)$ hacia la referencia, hasta alcanzar la banda de no arbitraje. Esto permite calcular los ticks ideales hasta la banda sin simulación. La simulación aborda las dos partes de la pregunta dentro de la grilla ensayada: cómo la capacidad de arbitraje y la profundidad de liquidez condicionan el seguimiento poschoque, y con qué precisión los mapas reducidos describen las trayectorias resultantes en horizonte finito. El ruido, la tendencia y las reglas de retiro hacen que estos resultados dependan de este micromundo con agentes heterogéneos.

El modelo basado en agentes es deliberadamente pequeño y reproducible. Representa una piscina de producto constante, una referencia exógena y reglas estilizadas de ruido, tendencia, arbitraje y proveedor de liquidez. La regla de retiro se activa en todas las ejecuciones de la condición con retiro por umbral; por eso se trata como una intervención fija y no como evidencia de respuestas heterogéneas de proveedores. En el diseño principal varía la profundidad y se mantienen fijos los tamaños absolutos de órdenes; un análisis separado los escala con las reservas. Por tanto, el experimento muestra cómo las conclusiones dependen de la razón entre el tamaño de orden y las reservas y de la regla de pánico, no cómo la profundidad determina por sí sola la respuesta de un mercado real.

La cláusula de modelos reducidos de la pregunta de investigación se evalúa como descripción predictiva del error de seguimiento continuo en trayectorias reservadas del simulador. Las etiquetas de ventana finita resumen esas trayectorias; no son regímenes observados independientemente ni objetivos de una prueba de clasificación. El mapa de dos estados no supera al autorregresivo escalar. La comparación del mapa sensible a la banda ofrece una mejora pequeña, concentrada en una condición sintética, y recibe como entrada la trayectoria futura realizada del CEX; es un diagnóstico de reproducción, no un pronóstico en tiempo real. Los resultados no establecen un clasificador de regímenes ni estabilidad global.

La implicación de ingeniería más útil es un resultado de medición específico del simulador. En E3, al ponderar por igual las nueve celdas del diseño, el 31.2\,\% de las trayectorias recuperadas salió después de la banda de no arbitraje; la tasa condicional fue 98.8\,\% con factor de profundidad 0.5 y cero con factores 1 y 2. Una medida posterior del error pico también muestra un aumento en E3 solo con la menor profundidad y órdenes fijas, aunque el retiro y el límite de arbitraje ligado a las reservas siguen confundidos. Estos patrones motivan registrar por separado la primera entrada, el retorno sostenido, las salidas posteriores y el error pico al final del tick en una futura reproducción sincronizada; no son tasas validadas ni requisitos operativos para mercados reales.

El documento se organiza así. El capítulo~\ref{cap:problema} delimita la pregunta de investigación y revisa antecedentes; el capítulo~\ref{cap:justificacion} explica el aporte científico y metodológico. Los capítulos~\ref{cap:hipotesis} y~\ref{cap:objetivos} presentan la proposición analítica, las preguntas de evaluación y los objetivos. El capítulo~\ref{cap:metodologia} especifica el modelo y el diseño de simulación; el capítulo~\ref{cap:resultados} informa los hallazgos analíticos y computacionales, y el capítulo~\ref{cap:discusion} discute su alcance. El capítulo~\ref{cap:futuro} expone el trabajo necesario para reproducir eventos empíricos y ampliar el modelo; la bibliografía y los anexos técnicos cierran el documento.

\chapter{Planteamiento del problema y estado del arte}
\label{cap:problema}

\section{Delimitación del problema}

Un CEX permite que los participantes introduzcan órdenes de compra y venta cuya ejecución depende de las órdenes disponibles en el libro. Un AMM de producto constante, en cambio, permite intercambios contra una piscina de dos activos. En la representación ideal de este último, el precio marginal se obtiene de la relación entre reservas, y la profundidad disponible determina la diferencia entre ese precio y el precio medio al que se ejecuta una transacción. Un cambio en el mercado centralizado no modifica directamente las reservas del AMM: para que su precio se ajuste deben producirse transacciones.

Se considera un par de activos $X/Y$ negociado en dos mecanismos. El CEX se representa por una serie de precios externos en unidades de $Y$ por unidad de $X$; el DEX, por un AMM. Esta simplificación permite controlar la perturbación aplicada a la referencia y observar solo la respuesta descentralizada. No se estudia un libro de órdenes completo ni una realimentación bidireccional entre mercados; la elección excluye de forma deliberada la profundidad del CEX y la sincronización de bloques, que exigirían otra infraestructura experimental.

El error de seguimiento es $e_t=\log(\pD_t/\pC_t)$. Su signo indica si el precio marginal del AMM está por encima o por debajo de la referencia. Los costos de transacción impiden interpretar toda diferencia como una oportunidad ejecutable, por lo que se distingue el error bruto del exceso fuera de una banda operativa. La duración de la discrepancia se mide sobre un horizonte explícito. Que el error vuelva a la banda no demuestra, por sí solo, estabilidad global de un sistema estocástico.

El problema tiene dos niveles. En el microscópico interesa cómo el arbitraje, las órdenes de los agentes y las decisiones de retiro alteran las reservas. En el reducido interesa si una relación de transición de pocas variables permite describir la respuesta del modelo completo. Distintas configuraciones pueden producir un error agregado similar; en tal caso, una descripción reducida no identifica necesariamente el mecanismo que lo generó. Por eso los experimentos intervienen sobre un mecanismo a la vez y conservan las demás condiciones.

\section{Sistemas complejos y mercados artificiales}

Los sistemas complejos adaptativos comprenden componentes heterogéneos que interactúan con su entorno y entre sí. Holland \cite{holland1995hidden} destaca la adaptación y la organización a partir de reglas locales; su obra se usa aquí solo como marco conceptual y no sustenta afirmaciones sobre mercados. En finanzas esta perspectiva es útil porque los precios observados modifican las decisiones posteriores, y la agregación de órdenes puede producir retroalimentaciones que no aparecen al estudiar una decisión aislada.

La economía computacional basada en agentes estudia procesos económicos mediante reglas ejecutables y poblaciones heterogéneas \cite{tesfatsion2003agent}, y su aplicación a las finanzas cuenta con revisiones consolidadas \cite{lebaron2006agent,hommes2006heterogeneous,axtell2025agent}. Un modelo basado en agentes (ABM) no exige que todos los participantes compartan una expectativa ni que el sistema alcance un equilibrio de inmediato. A diferencia de una representación agregada, conserva estados individuales y permite intervenciones contrafactuales sobre políticas particulares. Sin embargo, más detalle no equivale a más capacidad explicativa: con demasiadas reglas libres, múltiples configuraciones pueden explicar las mismas observaciones.

El mercado artificial de Palmer y colaboradores \cite{palmer1994artificial} y el modelo de Brock y Hommes \cite{brock1998heterogeneous} son antecedentes para estudiar dinámicas endógenas producidas por estrategias heterogéneas; Lux y Marchesi \cite{lux1999scaling} examinan regularidades financieras en un modelo estocástico de agentes. La implicación para este trabajo no es que un AMM reproduzca automáticamente esos resultados, sino que comparar políticas de ruido y de tendencia permite indagar cuánto de la dinámica depende de la heterogeneidad y cuánto de la estructura institucional. Además, los agentes de inteligencia cero \cite{gode1993allocative,farmer2005predictive} justifican que un flujo de órdenes aleatorio, sin información ni objetivo, sea una representación mínima legítima de los negociadores de ruido.

Los hechos estilizados de los mercados financieros, como las colas pesadas de los retornos y la persistencia de la volatilidad \cite{cont2001empirical}, proporcionan criterios de comparación. Reproducirlos de forma aislada es insuficiente como validación, porque distribuciones semejantes pueden emerger de mecanismos distintos, y un horizonte corto de simulación no permite estimar con precisión características de cola. Por eso el seguimiento de precios y la respuesta a una perturbación son aquí los resultados primarios, y los hechos estilizados quedan como diagnósticos secundarios sujetos a suficiencia de datos.

\section{Mercados artificiales de criptoactivos}

Cocco, Concas y Marchesi \cite{cocco2017artificial} trasladan el enfoque basado en agentes al estudio de un mercado artificial de Bitcoin, y Cocco, Tonelli y Marchesi \cite{cocco2019agent} lo extienden a la negociación de ese activo. Esos trabajos muestran que la investigación de criptoactivos mediante ABM antecede a la infraestructura DeFi actual. El mecanismo de interés aquí no es la emisión del activo ni la minería, sino la negociación contra reservas, y esos modelos no representan un AMM.

Comparar trabajos exige distinguir el mecanismo de formación de precios, el origen de las órdenes y el procedimiento de validación. Un mercado artificial que forma precios con un libro de órdenes y uno que usa una función de conservación no son intercambiables. Una propiedad matemática del AMM, como el aumento del impacto al disminuir sus reservas, no debe confundirse con una propiedad estadística producida por el comportamiento de los agentes. Este trabajo conserva ambos niveles y asigna a cada uno una prueba distinta.

\section{Intercambios descentralizados y AMM}

Schär \cite{schar2021defi} describe la infraestructura financiera basada en contratos inteligentes, y Xu y colaboradores \cite{xu2023sok} sistematizan los mecanismos de intercambio con AMM; Bartoletti, Chiang y Lluch-Lafuente \cite{bartoletti2022theory} ofrecen una teoría formal de los AMM. El caso de producto constante es suficientemente simple para derivar la relación entre el tamaño de una orden y el cambio de precio. Un AMM de liquidez concentrada, en cambio, distribuye posiciones por intervalos de precio y no puede representarse sin más con una única pareja de reservas globales.

El documento técnico de Uniswap v2 \cite{adams2020uniswap} permite distinguir el intercambio ideal sin comisión del intercambio con comisión retenida por la piscina. La condición $xy=k$ es una identidad durante un intercambio sin costos; al incorporar comisiones o entradas y salidas de proveedores de liquidez, el producto de reservas cambia. Una implementación debe registrar esos cambios de manera contable y no usar como invariante una igualdad que solo vale bajo supuestos específicos. Lehar y Parlour \cite{lehar2025uniswap} y Capponi y Jia \cite{capponi2025liquidity} analizan la provisión de liquidez y el comportamiento de los proveedores en DEX, lo que ubica la regla de retiro de este trabajo en un campo de estudio establecido.

Angeris y Chitra \cite{angeris2020cfmm} estudian los creadores de mercado de función constante y sus propiedades como oráculos de precio. Angeris y colaboradores \cite{angeris2020uniswap} analizan formalmente los mercados de producto constante, muestran que bajo ciertas condiciones siguen de cerca el precio de un mercado de referencia y lo verifican con una simulación basada en agentes a gran escala. El antecedente es directamente relevante: demuestra que el seguimiento de un precio externo por un AMM, con arbitrajistas y agentes heterogéneos, no es por sí mismo una brecha original. Su enfoque se usa como punto de partida para definir las pruebas de seguimiento y delimitar el uso del mapa reducido.

\section{Arbitraje, costos y liquidez adaptativa}

El arbitraje acerca el precio del AMM al precio externo cuando la discrepancia permite una operación rentable. Con comisión, el arbitraje solo es rentable si la discrepancia supera un umbral \cite{milionis2025fees}; con capital finito, aversión al riesgo o fricciones de liquidación, el ajuste es incompleto \cite{shleifer1997limits}. En criptoactivos, las desviaciones de precio entre mercados persisten por restricciones de capital y por los tiempos de liquidación \cite{makarov2020trading,hautsch2024building}, y el costo del arbitraje recae sobre los proveedores de liquidez, que pierden frente a los arbitrajistas cuando el precio de referencia se mueve \cite{milionis2022lvr}. La magnitud del ajuste depende del tamaño de la operación y de la capacidad disponible. Una regla que ajusta siempre el precio a su objetivo equivale a suponer capacidad suficiente y ejecución inmediata; introducir una capacidad máxima de negociación permite estudiar los casos en que el ajuste requiere varios pasos.

Campbell y colaboradores \cite{campbell2025optimal} analizan la elección de comisiones y la rentabilidad de proveedores pasivos en un modelo en que un AMM opera en paralelo con un CEX, con simulaciones y calibración con datos. Su objeto es la rentabilidad de la provisión y la determinación de comisiones, mientras que aquí se estudia el error de seguimiento y el alcance descriptivo de un mapa local. Esta distinción es una delimitación temática, no una prueba de que no existan otros trabajos sobre el mismo objeto. Conviene además señalar que el trabajo de Campbell y colaboradores es un preprint, igual que el de Bergault y colaboradores \cite{bergault2025exit} que se discute más adelante; Dean y colaboradores \cite{dean2022arbitrage} presentan un modelo de arbitraje con un AMM en un documento de trabajo que no ha sido arbitrado. Los tres se citan como antecedentes, no como resultados establecidos.

Palit \cite{palit2022shuffled} reordena aleatoriamente eventos históricos de Uniswap y los reproduce a través del mecanismo del AMM, lo que reduce el mercado real a inteligencia cero y permite medir su efecto en la pérdida impermanente, el deslizamiento y la eficiencia de precios. La lección metodológica es preservar un flujo y modificar una sola dimensión para aislar su efecto. Aquí se aplica la idea a pares de simulaciones con perturbaciones y flujos exógenos comunes que cambian una regla de arbitraje o de provisión de liquidez; no se afirma haber replicado el estudio histórico. Otros modelos basados en agentes sobre AMM son escasos: Sabaté-Vidales y Šiška \cite{sabatevidales2023variable} simulan comisiones variables en mercados de producto constante, y Cohen y colaboradores \cite{cohen2023inefficiency} combinan una perspectiva de cobertura con simulaciones multiagente y concluyen que, en promedio, las comisiones no compensan el riesgo de mercado. No se encontró un ABM previo que represente un retiro de liquidez dependiente del error de seguimiento.

Bergault, Bieber y Sánchez-Betancourt \cite{bergault2025exit} tratan la salida de un proveedor de liquidez como un problema de parada óptima, en el que el momento de retirarse depende del estado de la piscina y del precio de referencia. Estos antecedentes delimitan H2: no se propone descubrir que el retiro depende del estado, sino evaluar una regla heurística concreta y su efecto sobre la respuesta de precios bajo intervenciones controladas. La regla no reproduce la política óptima de esos autores y es una decisión de modelado propia, apoyada en ese antecedente y en el análisis de Lehar y Parlour \cite{lehar2025uniswap}; la analogía con las corridas bancarias \cite{diamond1983bank} se usa solo como intuición y no como evidencia sobre AMM.

La salida proporcional de reservas no cambia de inmediato el precio marginal, pero reduce el tamaño de la piscina respecto de una orden posterior. Esta distinción separa una propiedad mecánica verificable de una hipótesis sobre comportamiento adaptativo. El efecto mecánico se establece con una identidad; el adaptativo exige comparar trayectorias en que la decisión de retiro responde a observaciones previas.

Sobre la formación de precios entre CEX y DEX, Capponi, Jia y Yu \cite{capponi2026price} estudian el descubrimiento de precios en intercambios descentralizados, y Alexander y colaboradores \cite{alexander2025price} el descubrimiento de precios y la eficiencia de las piscinas de Uniswap. Esa literatura es empírica y apoya la pertinencia del tema; este trabajo no usa datos de mercado ni reproduce esos resultados.

\section{Modelos dinámicos reducidos}

El ABM puede representarse como una transformación discreta de su estado completo. Reducirlo a dos variables implica descartar información de inventarios, reservas y memoria de los agentes. El resultado no tiene necesariamente la propiedad de Markov: un mismo par de observables puede corresponder a estados internos distintos y, por tanto, a transiciones distintas. Por eso el mapa reducido se formula como aproximación local y no como identidad global. Ajustar mapas de grano grueso a la salida de simuladores microscópicos tiene antecedentes metodológicos \cite{gear2003equationfree}, y también el uso de sustitutos estadísticos de un ABM \cite{lamperti2018surrogates}.

La linealización alrededor de un punto de operación exige declarar el dominio y el significado de las variables \cite{roubal2010linearization}. Con una banda de no arbitraje, la función que convierte el error bruto en exceso no es diferenciable en los límites de la banda, de modo que conviene distinguir las regiones activa positiva, activa negativa e inactiva; combinar submodelos lineales locales por regiones tiene precedentes en la identificación de sistemas \cite{takagi1985fuzzy}. Un criterio de estabilidad de Schur \cite{schur1917potenzreihen,jury1962simplified,elaydi2005difference} aplicado a una matriz local permite estudiar una dinámica lineal autónoma, pero no establece estabilidad del sistema estocástico completo.

Comparar el modelo completo con su aproximación debe ir más allá del ajuste dentro de la muestra. Si los coeficientes se estiman con las mismas configuraciones en que se evalúa el error, el desempeño queda sobreestimado. Por eso el contraste se hace sobre configuraciones y semillas reservadas, con modelos base simples y métricas fijadas antes de la evaluación.

\section{Validación y reproducibilidad}

Fagiolo, Moneta y Windrum \cite{fagiolo2007critical} y Windrum, Fagiolo y Moneta \cite{windrum2007empirical} discuten las dificultades de validar empíricamente los ABM económicos. La verificación busca establecer si la implementación respeta la especificación; la validación considera si la especificación es adecuada para el fenómeno que representa \cite{sargent2013verification}. La distinción es indispensable en un modelo financiero: una identidad de conservación bien implementada no prueba que una política de negociación represente a participantes reales.

Grimm y colaboradores \cite{grimm2020odd} proponen el protocolo ODD para describir modelos basados en agentes. Su uso permite especificar entidades, variables, escalas, calendario de procesos, inicialización e información disponible (anexo~\ref{sec:odd}). El protocolo no reemplaza el diseño de experimentos; se complementa con un registro de escenarios, parámetros, semillas y resultados para que las comparaciones puedan repetirse \cite{railsback2019agent,thiele2014facilitating}.

\section{Síntesis comparativa y brecha delimitada}

La tabla~\ref{tab:antecedentes} resume lo que cada grupo de antecedentes aporta al problema y lo que implica para el estudio.

\begin{table}[htbp]
\centering
\caption{Comparación focalizada de antecedentes.}
\label{tab:antecedentes}
\small
\begin{tabularx}{\textwidth}{@{}>{\raggedright\arraybackslash}p{3.3cm} >{\raggedright\arraybackslash}X >{\raggedright\arraybackslash}X@{}}
\toprule
Antecedente & Aporte al problema & Implicación para el estudio \\
\midrule
Brock y Hommes \cite{brock1998heterogeneous}; Lux y Marchesi \cite{lux1999scaling} & Heterogeneidad y dinámica financiera. & Comparar reglas de ruido y de tendencia; no atribuir toda la dinámica a la emergencia. \\
Cocco et al. \cite{cocco2017artificial,cocco2019agent} & Mercados artificiales de Bitcoin. & Reconocer los antecedentes de ABM en criptoactivos; distinguir el mecanismo AMM. \\
Angeris et al. \cite{angeris2020uniswap} & Seguimiento de precios y simulación de agentes en Uniswap. & Usar la banda de arbitraje y reconocer la similitud con el sistema propuesto. \\
Xu et al. \cite{xu2023sok}; Bartoletti et al. \cite{bartoletti2022theory} & Taxonomía y teoría formal de los AMM. & Limitar el dominio al producto constante. \\
Palit \cite{palit2022shuffled} & Alteración controlada del orden de eventos históricos. & Separar efectos de reglas y de secuencias con flujos comunes. \\
Campbell et al. \cite{campbell2025optimal} & CEX/DEX, rentabilidad de proveedores y comisiones. & Distinguir el propósito del mapa reducido del problema de las tarifas óptimas. \\
Bergault et al. \cite{bergault2025exit} & Retiro de liquidez como parada óptima. & Evaluar una regla heurística, no la política óptima. \\
Gear et al. \cite{gear2003equationfree}; Roubal et al. \cite{roubal2010linearization} & Mapas de grano grueso y linealización local. & Declarar dominio de validez y comparar fuera de muestra. \\
Grimm et al. \cite{grimm2020odd}; Fagiolo et al. \cite{fagiolo2007critical} & Documentación y validación de ABM. & Separar verificación, experimentación y validación. \\
\bottomrule
\end{tabularx}
\end{table}

La revisión focalizada permite precisar la pregunta sin sostener una ausencia universal de modelos previos. El problema específico es establecer, para el micromundo definido, la relación entre restricciones de arbitraje, profundidad y respuesta adaptativa, y evaluar si un mapa discreto de bajo orden caracteriza la corrección o la persistencia del error en configuraciones distintas de las usadas para estimarlo. La aportación potencial es metodológica: una comparación explícita de representaciones, con condiciones de falsación y regiones de desacuerdo. La originalidad respecto de la literatura completa queda sujeta a una revisión bibliográfica ampliada.

\section{Pregunta de investigación}

\begin{quote}
\emph{¿Bajo qué condiciones de capacidad de arbitraje y profundidad de liquidez un AMM de producto constante sigue un precio CEX exógeno después de una perturbación, y con qué precisión un mapa dinámico reducido describe los regímenes de seguimiento obtenidos en un micromundo con agentes heterogéneos?}
\end{quote}

La pregunta no comprende la predicción del CEX ni pretende establecer causalidad sobre mercados reales. El universo de inferencia está formado por las configuraciones de un modelo computacional explícitamente especificado. Extender las conclusiones exigiría evidencia externa.

\chapter{Justificación}
\label{cap:justificacion}

\section{Justificación científica}

El seguimiento de un precio de referencia por un AMM ya se ha estudiado desde el análisis formal y la simulación \cite{angeris2020uniswap}, y también la rentabilidad de la provisión de liquidez frente a un CEX \cite{campbell2025optimal,milionis2022lvr}. Lo que este trabajo añade no es el fenómeno sino la forma de examinarlo: separa el efecto de la capacidad de arbitraje y el de la reacción de los proveedores mediante comparaciones pareadas con flujos de órdenes comunes, y contrasta una descripción reducida de dos variables con el simulador del que proviene.

La relevancia científica está en la delimitación de lo que cada representación puede afirmar. Una identidad analítica, como el aumento del impacto al disminuir las reservas, no requiere experimentación; una afirmación sobre arbitraje limitado o retiro adaptativo sí, porque depende de reglas de comportamiento y de un horizonte. Distinguir ambos niveles evita presentar como hallazgo empírico una consecuencia algebraica y evita también descartar como trivial un efecto que solo aparece bajo ciertas profundidades. El resultado del capítulo~\ref{cap:resultados} ilustra el punto: el efecto de la liquidez adaptativa sobre la respuesta máxima no es uniforme, sino que depende de la profundidad ensayada.

Además, la formulación como contraste confirmatorio, con hipótesis, métricas y criterios de rechazo fijados antes de simular las semillas de confirmación, responde a una crítica frecuente a los ABM: que se ajustan a posteriori hasta mostrar el fenómeno deseado \cite{fagiolo2007critical}. El protocolo no elimina esa crítica, pero hace explícito qué se decidió antes de observar los resultados y qué no. La señal de que el protocolo es útil es que también documenta resultados desfavorables a la expectativa inicial: el mapa reducido no supera a un autorregresivo escalar, y la precisión objetivo para H1 no se alcanzó con el número de semillas previsto.

\section{Justificación metodológica}

Una comparación sencilla entre un escenario con arbitraje y otro sin él no permite atribuir la diferencia al arbitraje si los flujos de órdenes difieren entre ambos. La comparación pareada, con los mismos choques exógenos y las mismas órdenes de ruido pregeneradas, aísla la regla modificada, y es la lógica que Palit \cite{palit2022shuffled} aplica al reproducir eventos históricos bajo un agente sin inteligencia. La inferencia con la semilla como unidad independiente evita tratar como muestras distintas las celdas de un mismo flujo, y el control de Holm limita el riesgo de falsos positivos al evaluar varias hipótesis a la vez.

El uso de configuraciones y semillas reservadas para evaluar el mapa reducido responde a un problema conocido en la identificación de modelos: el ajuste dentro de la muestra sobrestima la capacidad de descripción \cite{lamperti2018surrogates}. Ajustar el mapa con un conjunto de exploración y evaluarlo en otro permite distinguir entre una descripción que generaliza y una que solo reproduce las trayectorias de las que se estimó. La comparación contra la persistencia y contra un modelo autorregresivo escalar establece qué aporta el segundo estado del mapa. Sin esas referencias, un error pequeño no sería interpretable, porque la persistencia ya es una predicción competitiva a un paso.

Finalmente, el registro de cada ejecución con identificador, semilla, parámetros y verificación de invariantes \cite{grimm2020odd,railsback2019agent} permite que otra persona repita el experimento y compruebe las cifras publicadas. La reproducibilidad completa del experimento confirmatorio (identificador de ejecución, sumas de verificación y comando) se documenta en el anexo~\ref{cap:anexos}.

\section{Justificación aplicada y pedagógica}

Los resultados no están destinados a prescribir estrategias de inversión ni a calibrar un protocolo real. Su utilidad es de otro orden. Primero, para el diseño de mercados descentralizados, muestran que la profundidad inicial y la capacidad de arbitraje condicionan la rapidez de corrección, y que la reacción de los proveedores solo amplifica una perturbación bajo ciertas condiciones; esa información orienta qué parámetros conviene observar en un análisis posterior con datos. Segundo, en la formación de ingenieros de sistemas, el micromundo ofrece un caso en que la relación entre mecanismos locales, restricciones y comportamiento agregado puede explorarse con un código pequeño y verificable. Tercero, el sitio web asociado al trabajo reproduce el escenario con una interfaz interactiva, pensada para estudiar de manera intuitiva los efectos de reducir la profundidad o de limitar el arbitraje; su valor pedagógico no se ha evaluado con estudiantes, por lo que se menciona como línea de trabajo futuro (capítulo~\ref{cap:futuro}).

Los límites de estas aplicaciones están declarados: el modelo es sintético, no emplea datos de mercado y representa un único par con un AMM de producto constante. Cualquier uso fuera de ese dominio exige evidencia adicional.

\chapter{Proposición analítica y preguntas de evaluación}
\label{cap:hipotesis}

El estudio distingue una consecuencia algebraica de la regla de intercambio de dos preguntas que requieren simulación. Estas son afirmaciones sobre el micromundo especificado sin comisiones; no son hipótesis sobre mercados observados.

\section{H1: corrección ideal con arbitraje acotado}

Sea $e_t=\log(p_t^D/p_t^C)$ y $D_t=(|e_t|-\gamma)_+$. En una piscina aislada sin comisión, con referencia fija, sin otras órdenes y sin una restricción vinculante de inventario, la orden máxima $q=cx$ cambia el precio logarítmico en $\delta_c=2\log(1+c)$. Como el arbitrajista se detiene en la banda, el exceso de error evoluciona según
\[
  D_{t+1}=\max(0,D_t-\delta_c),\qquad
  N_\gamma=\left\lceil\frac{D_0}{2\log(1+c)}\right\rceil \quad(c>0).
\]
Para $c=0$ no hay corrección finita. Este es un referente mecánico derivado de la ecuación del intercambio de producto constante, no un hallazgo empírico. La simulación sirve para cuantificar las desviaciones respecto de ese referente causadas por otros agentes, el orden de eventos, el horizonte finito y las reglas de inventario especificadas.

\section{H2: retiro, escala de órdenes y respuesta de seguimiento}

La pregunta de evaluación que responden las corridas disponibles es más acotada:
% QUESTION_LOCK:h2:BEGIN
\emph{Dentro de los niveles de capacidad y profundidad ensayados, ¿cómo afecta al seguimiento posterior al choque la regla especificada de retiro único de liquidez, y cómo varía ese efecto cuando el tamaño de las órdenes escala con las reservas o se desactiva la conducta de pánico?}
% QUESTION_LOCK:h2:END
La regla E3 es un disparador de umbral único, no una política adaptativa estimada: retira el 45\,\% de las reservas en cuanto $|e_t|>0.02$. Los choques ejecutados superan este umbral en todas las corridas E3. El desplazamiento máximo original respecto al precio del AMM previo al choque se trata solo como diagnóstico; no mide el sobrepaso frente a la nueva referencia.

\section{H3: precisión descriptiva de los mapas reducidos}

La cláusula del mapa en la pregunta de investigación se operacionaliza como descripción predictiva del error de seguimiento continuo en trayectorias reservadas del simulador. Las etiquetas de trayectoria por ventana finita resumen esas trayectorias; no son regímenes observados independientemente ni objetivos de una prueba de clasificación. H3 pregunta:
% QUESTION_LOCK:h3:BEGIN
\emph{¿Con qué precisión predice el mapa de dos estados las trayectorias reservadas del error de seguimiento frente a la persistencia y un AR(1) escalar, y cómo se desempeña un mapa sensible a la banda frente al mapa afín en la reproducción recursiva?}
% QUESTION_LOCK:h3:END
La comparación se refiere a trayectorias generadas por el simulador. La reproducción híbrida recibe la trayectoria futura realizada del CEX, por lo que su resultado no evalúa un pronóstico en tiempo real. No se propone una hipótesis de clasificación de regímenes ni de estabilidad.

\chapter{Objetivos}
\label{cap:objetivos}

\section{Objetivo general}

Evaluar las condiciones de seguimiento de precios y amplificación de perturbaciones en un mercado CEX/DEX simplificado mediante experimentación computacional basada en agentes y el contraste de una aproximación dinámica discreta de bajo orden.

\section{Objetivos específicos}

\begin{enumerate}[label=\textbf{OE\arabic*.}, leftmargin=3em]
  \item Formular y verificar un simulador discreto de un AMM de producto constante con una referencia CEX exógena, agentes de ruido, tendencia, arbitraje y provisión de liquidez, y un registro reproducible de las ejecuciones.
  \item Caracterizar el efecto de la capacidad de arbitraje sobre el error de seguimiento y el tiempo de recuperación tras una perturbación de la referencia.
  \item Evaluar si la reducción adaptativa de liquidez amplifica la respuesta del precio del AMM y bajo qué profundidades.
  \item Estimar un mapa dinámico discreto de dos variables a partir de observables del simulador y contrastar su capacidad de predicción fuera de muestra frente a modelos de referencia.
\end{enumerate}

\section{Correspondencia entre objetivos, hipótesis y evidencia}

La tabla~\ref{tab:oe-h} relaciona cada objetivo con la hipótesis que lo contrasta, el método y la evidencia con la que se reporta. El objetivo OE1 es instrumental: no se contrasta con una hipótesis propia, pero habilita los demás, porque sin un simulador verificado ninguna comparación sería interpretable.

\begin{table}[htbp]
\centering
\caption{Correspondencia entre objetivos específicos, hipótesis, método y evidencia.}
\label{tab:oe-h}
\small
\begin{tabularx}{\textwidth}{@{}l l >{\raggedright\arraybackslash}X >{\raggedright\arraybackslash}X@{}}
\toprule
Objetivo & Hipótesis & Método & Evidencia \\
\midrule
OE1 & H1--H3 (instrumental) & Especificación, invariantes y registro de ejecuciones & Nueve verificaciones de invariantes; identificador y sumas de verificación (anexo~\ref{cap:anexos}) \\
OE2 & H1 & Diseño pareado con capacidades de arbitraje 0, baja, alta y máxima & Área de exceso, tasa de recuperación y tiempo de retorno (sección~\ref{sec:res-h1}) \\
OE3 & H2 & Diseño pareado estático frente a adaptativo en tres profundidades & Respuesta máxima, mecanismo y desglose por celda (sección~\ref{sec:res-h2}) \\
OE4 & H3 & Estimación con semillas de exploración; evaluación en celdas y semillas reservadas & Error de predicción frente a persistencia y autorregresivo escalar; clasificación de regímenes (sección~\ref{sec:res-h3}) \\
\bottomrule
\end{tabularx}
\end{table}

\chapter{Metodología}
\label{cap:metodologia}

\section{Enfoque}

El estudio combina una derivación analítica con experimentos computacionales. Se especifica un micromundo discreto y las corridas pareadas modifican reglas seleccionadas mientras reutilizan los mismos choques exógenos y órdenes de ruido. El diseño sigue la lógica de ciencia del diseño \cite{hevner2004design} en cuanto a que se construyen el simulador y los mapas reducidos, se verifican y se evalúan con criterios explícitos, y se apoya en la metodología de modelado y simulación descrita en la literatura en español \cite{hernandez2014metodologia}. Las corridas de exploración estiman mapas candidatos; se usan semillas y celdas disjuntas para las comparaciones reservadas. Los resultados de simulación están condicionados al modelo codificado y no son inferencias de una muestra de mercados.

El procedimiento tiene cuatro pasos: (i) especificar el AMM, la referencia CEX y las reglas de los agentes; (ii) verificar la implementación con invariantes; (iii) estimar el mapa reducido con semillas y celdas de exploración; (iv) comparar de forma pareada sus resultados en semillas y celdas reservadas.

\section{Revisión bibliográfica focalizada}

La literatura se seleccionó por relevancia para cinco componentes del diseño: modelos basados en agentes en finanzas, mecánica de los AMM, arbitraje con costos y límites, modelos reducidos, y validación y reproducibilidad. Cada referencia se contrastó con el registro del editor o con un repositorio de preprints, y se anotaron las discrepancias de metadatos y el estado de revisión por pares (anexo~\ref{sec:protocolo-literatura}). Los preprints y documentos de trabajo se citan como tales. Los criterios de inclusión fueron la pertinencia directa con el mecanismo estudiado y la verificabilidad; se descartaron trabajos cuya relación con el problema era solo temática.

\section{Mecánica del AMM de producto constante}
\label{sec:amm}

El AMM tiene reservas $(x_t,y_t)$ de los activos $X$ e $Y$, con precio marginal
\begin{equation}
  \pD_t = \frac{y_t}{x_t},
  \label{eq:precio}
\end{equation}
y producto $k_t=x_ty_t$. En un intercambio sin comisión, la venta de una cantidad $q>0$ de $X$ al AMM conserva el producto y deja las reservas
\begin{equation}
  x' = x+q,\qquad y' = \frac{xy}{x+q}.
  \label{eq:reservas}
\end{equation}
El precio marginal posterior y el impacto logarítmico son
\begin{equation}
  p' = p\left(1+\frac{q}{x}\right)^{-2},\qquad I(q,x)=\left|\log\frac{p'}{p}\right| = 2\log\left(1+\frac{q}{x}\right).
  \label{eq:impacto}
\end{equation}
El deslizamiento, definido como la diferencia relativa entre el precio medio de ejecución y el precio marginal inicial, es $q/(x+q)$. Con una comisión proporcional $f$ retenida por la piscina, la cantidad efectiva que entra al invariante es $(1-f)q$, de modo que
\begin{equation}
  x'=x+q,\qquad y'=\frac{xy}{x+(1-f)q},\qquad k'\ge k.
  \label{eq:comision}
\end{equation}
La ecuación~\eqref{eq:comision} se incluye para precisar el supuesto del experimento: \textbf{el simulador ejecutado implementa únicamente intercambios sin comisión} ($f=0$). Los niveles de comisión de la tabla~\ref{tab:factores} son candidatos de diseño para trabajos posteriores y no se ejecutaron. La banda de no arbitraje sí se mantiene, pero como parámetro de tolerancia y no como consecuencia de una comisión.

Un retiro proporcional de liquidez, que multiplica ambas reservas por la misma fracción, mantiene el precio marginal y reduce $k$ en el cuadrado de esa fracción. Eso hace la piscina más sensible a una orden posterior de un tamaño dado; la identidad
\begin{equation}
  \frac{\partial I}{\partial x} = -\frac{2q}{x(x+q)} <0
  \label{eq:dI}
\end{equation}
muestra que el impacto decrece con la reserva. Es una propiedad mecánica y no una hipótesis.

\section{Referencia centralizada}

El precio CEX es exógeno y sigue
\begin{equation}
  \log \pC_{t+1} = \log \pC_t + \mu + \sigma\,\varepsilon_{t+1} + s_{t+1},
  \label{eq:cex}
\end{equation}
con deriva $\mu$, volatilidad $\sigma$, ruido gaussiano $\varepsilon$ y un choque determinista $s$ aplicado una sola vez en el tick 1. En todos los escenarios ejecutados $\mu=\sigma=0$, de modo que la única perturbación de la referencia es el choque. Esta elección es deliberada: aísla la respuesta del AMM a una perturbación conocida y evita confundirla con el ruido de la referencia. La volatilidad externa aparece en la tabla~\ref{tab:factores} como candidato de diseño, no ejecutado.

\section{Error de seguimiento, banda y arbitraje}

El error logarítmico es $e_t=\log(\pD_t/\pC_t)$. Se define una banda de no arbitraje de semiancho $\gamma$ (valor ejecutado $0.003$) y el exceso de error con signo
\begin{equation}
  d_t = \sgn(e_t)\,\bigl(|e_t|-\gamma\bigr)_+ ,\qquad D_t=|d_t|.
  \label{eq:exceso}
\end{equation}
Si $e_t>\gamma$, un arbitrajista vende $X$ al AMM hasta llevar el precio al borde de la banda. El tamaño que lo logra es
\begin{equation}
  q^{\ast} = x\left(\sqrt{\frac{\pD}{\pC e^{\gamma}}}-1\right),
  \label{eq:qstar}
\end{equation}
y la orden ejecutada es $\min\{q^{\ast},\,c\,x\}$, además limitada por el inventario del arbitrajista; $c$ es la capacidad de arbitraje, expresada como fracción máxima de la reserva negociable por tick. El caso $e_t<-\gamma$ es simétrico, con compra de $X$ pagada en $Y$. Con $c=0$ no se ejecuta ninguna orden. La capacidad es el único parámetro que gradúa el arbitraje en el diseño: $c=0.005$ es capacidad baja, $0.03$ alta y $0.10$ máxima.

\section{Agentes}

El estado contable compacto registrado en el tick $t$ es
\begin{equation}
  S_t = \bigl(x_t,\,y_t,\,\pC_t,\,W_t\bigr),
  \label{eq:estado}
\end{equation}
donde $W_t$ reúne los inventarios de los agentes. Este vector es un estado contable compacto; no se afirma que $S_t$ por sí solo sea un estado completo de Markov. La regla de tendencia depende del cambio de precio del AMM respecto del tick anterior; por tanto, el estado de transición debe conservar ese precio o una memoria equivalente del agente. También deben conservarse las banderas de reglas de una sola activación que no puedan inferirse de los inventarios. En una formulación estocástica, la siguiente orden aleatoria puede representarse como entrada explícita; la reproducción determinista requiere además el estado del generador o el índice de la secuencia de órdenes pregenerada. El presente análisis no formaliza ese estado aumentado del simulador completo. Cada agente opera con inventario finito; los agentes de ruido y tendencia disponen de dos veces las reservas iniciales y el arbitrajista de cinco veces. En las corridas Monte Carlo archivadas ninguna orden estuvo limitada por capital (0 de 13\,500).

\paragraph{Ruido.} Con probabilidad $0.7$ por tick, emite una orden de lado aleatorio por el $0.4\,\%$ de la reserva inicial del activo de entrada. El calendario de órdenes se genera antes de la simulación a partir de la semilla, de modo que los diseños pareados comparten el mismo flujo.

\paragraph{Tendencia y pánico.} Tras la orden de ruido, calcula el retorno del precio del AMM respecto del final del tick anterior. Si es menor o igual a $-1.5\,\%$, vende una vez por ejecución el $2\,\%$ de la reserva inicial de $X$ (venta de pánico). En otro caso, si su magnitud alcanza $0.1\,\%$, sigue la tendencia con una orden del $0.6\,\%$ de la reserva inicial del activo de entrada.

\paragraph{Arbitraje.} Ejecuta la regla de la sección anterior con capacidad $c$.

\paragraph{Provisión de liquidez.} Con la regla de retiro por umbral, un proveedor retira una sola vez, de forma proporcional, el $45\,\%$ de las reservas (conserva el $55\,\%$) si $|e_t|>0.02$. En el escenario estático no retira. Esta regla es una decisión de modelado propia: no reproduce la política de parada óptima de la literatura \cite{bergault2025exit} ni se calibró con datos.

\section{Calendario de un tick}

En cada tick: (1) se aplica el choque de la referencia (solo en el tick 1); (2) se ejecuta la orden de ruido pregenerada; (3) el agente de tendencia o pánico evalúa el retorno del AMM; (4) el proveedor aplica el umbral de retiro único; (5) el arbitrajista ejecuta su orden acotada; (6) se registran el estado y las métricas y se comprueban las identidades de conservación de activos (tolerancia $10^{-8}$ veces el mayor total del sistema) y la positividad y finitud de las reservas. El horizonte es de 20 ticks, con reservas iniciales $x_0=y_0=10\,000$ y precio inicial 1. La profundidad se modifica con un factor $\kappa$ que multiplica las reservas y divide las fracciones de las órdenes, de modo que el tamaño absoluto de las órdenes de ruido se mantiene (40 unidades) y el impacto de una venta de ruido es $0.016$, $0.008$ y $0.004$ con $\kappa=0.5,1,2$.

\section{Mapa dinámico de bajo orden}

Las dos variables del mapa son el exceso de error $d_t$ y la presión neta de compra $z_t$, definida como la suma con signo de las órdenes de ruido y tendencia (se excluye el arbitraje) dividida por la reserva inicial de $Y$. Con la entrada exógena $u_t=-\Delta\log\pC_t$, el mapa es
\begin{align}
  d_{t+1} &= A_{11}\,d_t + A_{12}\,z_t + u_{t+1}, \label{eq:mapa-d}\\
  z_{t+1} &= A_{21}\,d_t + A_{22}\,z_t. \label{eq:mapa-z}
\end{align}
Las entradas $A_{ij}$ se ajustan con mínimos cuadrados sin intercepto: la primera fila usa el objetivo $d_{t+1}-u_{t+1}$ y la segunda el objetivo $z_{t+1}$. Esta notación distingue los coeficientes ajustados del parámetro $c$ de capacidad de arbitraje. Como $d_t$ vale cero dentro de la banda y cambia de forma en su límite, un único mapa afín no puede representar exactamente ese umbral. Si el predictor $d_t$ es numéricamente nulo (caso sin variación del exceso), se usa solo el segundo predictor y se marca la ecuación como degenerada. Los parámetros no se suponen universales: se estiman por condición.


\subsection{Estimación y evaluación reservada}

Cada condición tiene su propio ajuste, estimado solo con las ejecuciones de exploración (semillas 0 a 299, profundidad 1, choques $\pm0.1$). La evaluación usa el error absoluto medio a un paso sobre 20 transiciones de las ejecuciones reservadas y disjuntas. Los modelos de referencia son la persistencia, $\hat d_{t+1}=d_t+u_{t+1}$, y un autorregresivo escalar $d_{t+1}-u_{t+1}=a_s d_t$, ajustado por condición con los datos de exploración. Cada trayectoria recibe además una etiqueta descriptiva: \emph{en banda} si $|d_t|<10^{-9}$ en los ticks 1--5 y 16--20; \emph{contractiva} si no está en banda y la media de $|d_t|$ es menor en los ticks 16--20 que en los ticks 1--5; \emph{persistente} en otro caso. Estas etiquetas resumen trayectorias del simulador, no regímenes observados de forma independiente. Por tanto, la cláusula del mapa en la pregunta de investigación se evalúa como predicción del error continuo en condiciones reservadas; esta versión no estima la precisión de una clasificación de regímenes.

\subsection{Comparación con un mapa híbrido sensible a la banda}
\label{sec:hybrid-method}

Un análisis separado (modelo 0.5.0, análisis 0.1.0) compara el mapa afín L0 con el candidato híbrido Hyb. Esta comparación es distinta de H1, la hipótesis de la tesis sobre el arbitraje. Usa el error de seguimiento con signo $e_t=\log(p_t^D/p_t^C)$ y la misma variable de flujo normalizada $z_t$. La entrada externa conocida es $u_{t+1}=-\Delta\log p^C_{t+1}$. Se definen el exceso con signo $\phi_\gamma(e)=\operatorname{sgn}(e)\max(|e|-\gamma,0)$ y el indicador $I_t=\mathbf{1}(|e_t|>\gamma)$, con banda configurada $\gamma=0.003$. Las ecuaciones son
\begin{align}
  \text{L0:}\quad e_{t+1}&=e_t+u_{t+1}+A_{11} e_t+A_{12} z_t,& z_{t+1}&=A_{21} e_t+A_{22} z_t,\\
  \text{Hyb:}\quad e_{t+1}&=e_t+u_{t+1}+A_{11} e_t+A_{12} z_t+g_e\phi_\gamma(e_t)+h_e z_t I_t,&
  z_{t+1}&=A_{21} e_t+A_{22} z_t+g_z\phi_\gamma(e_t)+h_z z_t I_t.
\end{align}
Los coeficientes se ajustan sin intercepto y por separado para las cuatro condiciones operativas compartidas. Por tanto, Hyb añade dos términos dependientes de la banda por variable predicha; no añade ramas de estado para la capacidad o el retiro. La cuadrícula de selección usa las semillas 3000--3299, profundidad 1 y choques $\pm0.10$. La evaluación intacta usa semillas 4000--4299, profundidades 0.5 y 2, y choques $\pm0.05$ y $\pm0.20$. La métrica primaria es el MAE recursivo por semilla en $e$, promediado sobre las 32 celdas reservadas por semilla y 20 transiciones. La predicción recursiva recibe la trayectoria exógena conocida del CEX y alimenta sus propios valores predichos $(e,z)$. El intervalo bootstrap percentílico pareado usa 10 000 remuestreos de las 300 unidades semilla; los contrastes por celda y las demás métricas de estado son resultados secundarios descriptivos. La evaluación compara errores de mapas reducidos en trayectorias generadas por el simulador; no prueba pronósticos en tiempo real ni la estabilidad del sistema híbrido completo.
\ifdefined\blindreviewmode
Los archivos y metadatos específicos de esta ejecución se omiten en la copia anónima.
\else
La configuración exacta, las sumas de verificación del código y los resultados se registran en \code{experiments/cex\_dex\_mvp/hybrid\_config.json}.
\fi

La ecuación de $z_{t+1}$ no recibe como entrada el flujo de órdenes estocástico realizado en el tick siguiente. Su objetivo incluye sorteos futuros de ruido o tendencia que este mapa basado solo en el estado no observa; por tanto, los errores de presión de flujo diagnostican una especificación incompleta de las entradas y no que el flujo del simulador sea determinista.

\section{Diseño experimental}

\begin{table}[htbp]
\centering
\caption{Escenarios experimentales. Las capacidades son fracciones de la reserva negociable por tick.}
\label{tab:escenarios}
\small
\begin{tabularx}{\textwidth}{@{}l >{\raggedright\arraybackslash}X l l l@{}}
\toprule
Escenario & Descripción & Capacidad & Liquidez & Región \\
\midrule
E0 & Referencia sin arbitraje & 0 & estática & primaria \\
E1-baja & Arbitraje de capacidad baja & 0.005 & estática & primaria \\
E1-alta & Arbitraje de capacidad alta & 0.03 & estática & primaria \\
E2 & Reducción de profundidad (factor $\kappa$) & --- & --- & control analítico \\
E3 & Arbitraje alto con retiro por umbral & 0.03 & por umbral & primaria \\
E1-máx & Arbitraje de capacidad máxima & 0.10 & estática & extrapolación \\
\bottomrule
\end{tabularx}
\end{table}

E2 no se simula como condición independiente: la profundidad entra como factor $\kappa$ en todos los escenarios y su efecto mecánico se presenta como tabla analítica (capítulo~\ref{cap:resultados}). La tabla~\ref{tab:factores} distingue los factores ejecutados de los candidatos de diseño.

\begin{table}[htbp]
\centering
\caption{Factores del diseño. Solo los niveles marcados como ejecutados forman parte de los resultados.}
\label{tab:factores}
\scriptsize
\setlength{\tabcolsep}{3pt}
\renewcommand{\arraystretch}{0.9}
\begin{tabularx}{\textwidth}{@{}>{\raggedright\arraybackslash}p{0.19\textwidth} >{\raggedright\arraybackslash}X >{\raggedright\arraybackslash}p{0.19\textwidth}@{}}
\toprule
Factor & Niveles & Estado \\
\midrule
Capacidad de arbitraje $c$ & 0; 0.005; 0.03; 0.10 & ejecutado \\
Profundidad $\kappa$ & 0.5; 1; 2 (reservado); 1 (exploración) & ejecutado \\
Choque de la referencia & $-0.05$; $-0.10$; $-0.20$ (reservado); $\pm0.10$ (exploración) & ejecutado \\
Retiro de liquidez & estático; por umbral & ejecutado \\
Comisión $f$ & niveles por definir & candidato, no ejecutado \\
Latencia del arbitraje & niveles por definir & candidato, no ejecutado \\
Volatilidad externa $\sigma$ & niveles por definir & candidato, no ejecutado \\
\bottomrule
\end{tabularx}
\end{table}

\subsection{Ajuste de mapas y comparaciones reservadas}

La exploración usa las semillas 0 a 299, profundidad 1 y choques $\pm0.10$. Las corridas reservadas usan las semillas 1000 a 1299, profundidades 0.5, 1 y 2 y choques $-0.05$, $-0.10$ y $-0.20$, con 5 condiciones y 9 celdas por semilla, o 13\,500 ejecuciones (2700 por condición). Los rangos de semillas son disjuntos. Para la comparación del mapa a un paso se evalúan las 8 celdas cuyo par (profundidad, choque) no pertenece a la grilla de exploración; se excluye la celda $(1.0,-0.10)$. Estas decisiones se fijaron antes de ejecutar el conjunto reservado y se registran en el anexo. Definen comparaciones Monte Carlo dentro de este modelo, no pruebas preregistradas sobre el comportamiento del mercado.

\section{Métricas}

\begin{table}[htbp]
\centering
\caption{Métricas del experimento.}
\label{tab:metricas}
\small
\begin{tabularx}{\textwidth}{@{}l >{\raggedright\arraybackslash}X l@{}}
\toprule
Métrica & Definición & Uso \\
\midrule
Área de exceso $A$ & Suma del exceso absoluto en 20 ticks & H1 descriptivo; comparar con referente ideal \\
Desplazamiento legado $M$ & Cambio máximo del log-precio AMM desde antes del choque & Diagnóstico H2; no es sobrepaso \\
Pico de horizonte completo $P$ & $\max_{0\le t\le20}|e_t|$ al final de cada tick & Diagnóstico legado dominado por el choque \\
Área posrecuperación $A_{\mathrm{post}}$ & $\sum_{t>t_c}(|e_t|-\gamma)_+$ & Resultado H2 exploratorio principal \\
Pico posrecuperación $P_{\mathrm{post}}$ & $\max_{t>t_c}(|e_t|-\gamma)_+$ & Resultado H2 exploratorio secundario \\
Episodios de salida & Observaciones consecutivas al final del tick fuera de banda tras $t_c$ & Diagnóstico descriptivo del mecanismo \\
Error de predicción & MAE a un paso sobre 20 transiciones & Comparación reservada H3 \\
Retorno sostenido & Dos ticks consecutivos con $|e_t|\le\gamma$ & Secundaria \\
Tiempo de retorno & Tick de inicio; 21 si no hay retorno sostenido & Secundaria \\
Profundidad mínima & $\min_t\sqrt{x_ty_t}$ & Secundaria \\
Activación del retiro & Si se aplicó el disparador fijo & Mecanismo \\
Órdenes limitadas por capital & Conteo & Verificación \\
Etiqueta de trayectoria & En banda, contractiva o persistente & Solo descriptiva \\
\bottomrule
\end{tabularx}
\end{table}

El tiempo de retorno está censurado a la derecha: una ejecución sin retorno sostenido recibe el valor 21 y se informa la media restringida, que no debe leerse como un tiempo esperado de recuperación.

El pico de horizonte completo usa solo estados registrados al final de cada tick y omite excursiones de precio entre actualizaciones. Incluye el choque impuesto y se conserva como diagnóstico legado, no como resultado H2 primario. Para las métricas posrecuperación, $t_c$ es la segunda observación del primer par de estados consecutivos registrados al final del tick que cumplen $|e_t|\le\gamma$; la suma y el máximo comienzan en el tick siguiente. Un episodio de salida es un tramo consecutivo de observaciones al final del tick fuera de banda tras $t_c$. La coincidencia con un tick de orden de pánico es descriptiva y no identifica causalidad intratick.

\subsection{Repetición exploratoria de H2}
\label{sec:h2-replay-method}

Tras la revisión crítica, el protocolo 0.7.0 repitió las semillas 2000--2299 de la ejecución de ingeniería archivada. Incluye 5\,400 trayectorias con órdenes fijas y tres profundidades/choques, 1\,800 trayectorias poco profundas con la orden de pánico desactivada y 1\,800 poco profundas con el límite de entrada de arbitraje fijado al valor calculado sobre las reservas iniciales; se ejecutaron E1-alta estático y E3 con retiro por umbral en cada variante, para 9\,000 trayectorias. La variante sin pánico conserva las órdenes de seguimiento de tendencia. La variante de capacidad fija usa $\min(c\,x_{\mathrm{inicial}},x_{\mathrm{activo}})$ como entrada máxima, con la reserva del activo de entrada pertinente. La repetición con límite original reprodujo exactamente las métricas archivadas de las 7\,200 trayectorias con órdenes fijas y sin pánico.

El contraste exploratorio principal es E3 menos E1-alta del área de exceso posrecuperación. Los contrastes por celda se condicionan a que ambas trayectorias se recuperen y tengan al menos un tick posconfirmación; se aplican 10\,000 remuestreos bootstrap por semilla y se informan la cobertura de recuperación y los tamaños pareados. Las proporciones de salida usan en todas las tablas el mismo denominador de trayectorias recuperadas con seguimiento posconfirmación y emplean intervalos de Wilson. Como se excluyen las trayectorias que no se recuperan, y es plausible que tengan mayor error residual, este $\Delta A_{\mathrm{post}}$ condicionado puede subestimar un estimando de fragilidad sobre todas las trayectorias. Estos análisis reutilizan flujos de semillas existentes, se especificaron después de examinar resultados previos y no constituyen un experimento independiente ni confirmatorio.

El diagnóstico posterior de razones de capacidad 0.8.0 repitió 1\,800 trayectorias E3 poco profundas con las mismas semillas: 900 con el límite fijado a las reservas iniciales y el pánico activo para verificar el control existente, y 900 con el mismo límite y la orden de pánico desactivada. La variante con pánico reprodujo 123, 101 y 53 episodios únicos de salida para los choques $-0.05$, $-0.10$ y $-0.20$, respectivamente; todos comenzaron en el tick del pánico. La variante sin pánico no tuvo salidas posteriores en sus 900 trayectorias. Es una comprobación posterior y dirigida del diagnóstico $R$, no un experimento independiente ni confirmatorio.

\section{Resúmenes Monte Carlo}
\label{sec:analisis-estadistico}

La semilla es la unidad de réplica Monte Carlo. Las celdas se promedian dentro de cada semilla y se informan las diferencias pareadas media y mediana con intervalos bootstrap percentílicos a partir de 10\,000 remuestreos \cite{efron1979bootstrap}. Estos intervalos caracterizan la variabilidad del proceso de órdenes aleatorio programado bajo reglas y niveles fijos; no representan incertidumbre entre mercados reales. Las pruebas previas de inversión de signos con 20\,000 sorteos devolvieron el mínimo valor resoluble para los tres contrastes informados; por ello, el ajuste de Holm produjo el mismo valor en el piso para cada uno. Esos valores $p$ y los objetivos previos de tamaño de muestra se omiten de la interpretación porque no añaden información sobre la magnitud del efecto ni el alcance del modelo. Los productos archivados siguen disponibles para reproducibilidad.

\section{Verificación, validación y trazabilidad}

La verificación comprende nueve comprobaciones automáticas (anexo~\ref{sec:verificacion}), que combinan identidades mecánicas y límites con comprobaciones de las reglas configuradas y la consistencia de las corridas pareadas. Algunas repiten comportamientos esperados que ya están incorporados en las reglas; aprobarlas no es respaldo independiente de afirmaciones causales.

La trazabilidad se apoya en un contrato de ejecución que registra identificador, versión del código, escenario, semilla, parámetros, series, métricas y sumas de verificación (anexo~\ref{sec:contrato}). El código y los datos de las figuras se documentan junto con los materiales del estudio \cite{calderon2026repository}. Los resultados completos no se versionan: se citan por identificador, sumas de verificación y el comando que los regenera.

\ifdefined\blindreviewmode
Los comandos, rutas e identificadores de ejecución de la comparación híbrida se omiten en esta copia para evitar que permitan localizar el repositorio.
\else
La comparación híbrida se reproduce con \code{pnpm experiment:mvp:hybrid}; sus exportaciones compactas para la web y la tesis se regeneran con \code{pnpm experiment:mvp:hybrid:export}. La ejecución \code{\HybridRunID} registra las configuraciones exactas del modelo y el análisis, los rangos de semillas, la versión de Python, los hashes de las fuentes y las sumas de verificación de los resultados.
\fi

\chapter{Resultados}
\label{cap:resultados}

\section{Alcance de los resultados}

Este capítulo informa identidades analíticas, comparaciones Monte Carlo sobre 13\,500 ejecuciones, una sensibilidad de ingeniería con 40\,500 trayectorias (sección~\ref{sec:engineering-sensitivity}) y una comparación reservada de mapas reducidos afines e híbridos sobre 9\,600 trayectorias (sección~\ref{sec:hybrid-results}). Todas las trayectorias provienen del micromundo especificado sin comisiones. Las corridas y los detalles de reproducción se documentan en el anexo; los identificadores se omiten en la copia anónima.

\section{Resultados analíticos: profundidad e impacto}
\label{sec:res-analitico}

Para una venta fija de $q=500$ unidades de $X$, el impacto $I(q,\kappa x)$ de la ecuación~\eqref{eq:impacto} crece al reducir la reserva (tabla~\ref{tab:impacto}). Con $x=10\,000$, reducir la reserva a la cuarta parte multiplica el impacto por 3.7 y el deslizamiento por 3.5.

\begin{table}[htbp]
\centering
\caption{Impacto de una venta de $q=500$ unidades según la reserva (resultado analítico).}
\label{tab:impacto}
\small
\begin{tabular}{@{}rrrr@{}}
\toprule
Factor de reserva & Reserva $x$ & Impacto $I$ & Deslizamiento (\%) \\
\midrule
1.00 & 10\,000 & 0.09758 & 4.76 \\
0.75 & 7500 & 0.12908 & 6.25 \\
0.50 & 5000 & 0.19062 & 9.09 \\
0.25 & 2500 & 0.36464 & 16.67 \\
\bottomrule
\end{tabular}
\end{table}

La figura~\ref{fig:impacto} muestra la curva completa. Con $I(q,\kappa x)=2\log(1+q/(\kappa x))$ y $\kappa<1$ se cumple $I(q,\kappa x)>I(q,x)$ para todo $q>0$, y la derivada de la ecuación~\eqref{eq:dI} es negativa. Este resultado es una consecuencia algebraica de la regla de producto constante y también sirve como control de la implementación (verificación 8 del anexo).

\subsection{Tasa ideal de corrección con arbitraje acotado}

En ausencia de otras órdenes y con referencia fija, un intercambio limitado a $q=cx$ cambia el precio logarítmico en una magnitud $\delta_c=2\log(1+c)$. Si el error está fuera de la banda, el exceso restante sigue $D_{t+1}=\max(0,D_t-\delta_c)$, y el número ideal de ticks hasta la banda es
\[
  N_\gamma=\left\lceil\frac{D_0}{2\log(1+c)}\right\rceil,\qquad D_0=(|e_0|-\gamma)_+.
\]
Para un choque logarítmico de magnitud $0.20$ y $\gamma=0.003$, se requieren idealmente 20 ticks con $c=0.005$, cuatro con $c=0.03$ y dos con $c=0.10$; con $c=0$ no hay corrección. Por tanto, el resultado de capacidad baja queda en gran medida implicado por la regla y el horizonte de 20 ticks. La simulación aporta información sobre cómo el ruido, la tendencia, el retiro y el orden de eventos alteran ese referente.

\begin{table}[htbp]
\centering
\caption{Corrección idealizada con referencia fija y sin otras órdenes.}
\label{tab:capacity-benchmark}
\small
\begin{tabular}{@{}rrr@{}}
\toprule
Capacidad $c$ & Corrección logarítmica por tick $2\log(1+c)$ & Ticks ideales hasta la banda \\
\midrule
0 & 0 & No hay corrección finita \\
0.005 & 0.00998 & 20 \\
0.03 & 0.05912 & 4 \\
0.10 & 0.19062 & 2 \\
\bottomrule
\end{tabular}
\end{table}

\begin{figure}[htbp]
\centering
\begin{tikzpicture}
\begin{axis}[width=0.78\textwidth,height=6cm,xlabel={Factor de reserva $\kappa$},ylabel={Impacto logarítmico $I$},grid=major,xmin=0.1,xmax=1.05]
\addplot[thick,blue!70!black] table[col sep=comma,x=reserve_factor,y=log_impact] {\datadir/impact_curve.csv};
\end{axis}
\end{tikzpicture}
\caption{Impacto logarítmico de una venta de $q=500$ unidades según el factor de reserva.}
\label{fig:impacto}
\end{figure}

\section{Panorama de las condiciones}

La tabla~\ref{tab:resumen} resume las cinco condiciones, con 2700 ejecuciones cada una (300 semillas por 9 celdas). E1-máx es una extrapolación fuera de los niveles principales de capacidad. Las 300 semillas son una decisión de réplica computacional, no el tamaño de una muestra de mercado.

\begin{table}[htbp]
\centering
\caption{Resumen de las condiciones reservadas (2700 ejecuciones por condición).}
\label{tab:resumen}
\small
\begin{tabular}{@{}lrrrrrr@{}}
\toprule
Condición & Área $A$ & $M$ legado & Recup. & Retorno\,$^{a}$ & Prof. mín. & Retiro \\
\midrule
E0 & 2.363 & 0.1012 & 4.6\,\% & 20.47 & 11667 & 0\,\% \\
E1-baja & 1.056 & 0.1305 & 56.8\,\% & 14.82 & 11667 & 0\,\% \\
E1-alta & 0.100 & 0.1220 & 100.0\,\% & 2.54 & 11667 & 0\,\% \\
E3 & 0.180 & 0.1335 & 96.6\,\% & 4.27 & 6417 & 100\,\% \\
E1-máx & 0.005 & 0.1197 & 100.0\,\% & 1.22 & 11667 & 0\,\% \\
\bottomrule
\multicolumn{7}{@{}l}{\footnotesize $^{a}$ Media restringida con censura en 21 ticks.}
\end{tabular}
\end{table}

El error pico medio es $0.176$ en E0, $0.127$ en E1-baja, $0.064$ en E1-alta, $0.073$ en E3 y $0.008$ en E1-máx. Ninguna orden estuvo limitada por capital en las 13\,500 ejecuciones.

\begin{figure}[htbp]
\centering
\begin{tikzpicture}
\begin{axis}[width=0.9\textwidth,height=6.4cm,xlabel={Tick},ylabel={$|d_t|$ medio},grid=major,legend pos=north east,legend style={font=\scriptsize},xmin=0,xmax=20,ymin=0]
\addplot[thick,red!70!black] table[col sep=comma,x=tick,y=mean_abs_excess,restrict expr to domain={\thisrow{series_id}}{10:10}] {\datadir/trajectories.csv};
\addplot[thick,orange!85!black] table[col sep=comma,x=tick,y=mean_abs_excess,restrict expr to domain={\thisrow{series_id}}{20:20}] {\datadir/trajectories.csv};
\addplot[thick,blue!70!black] table[col sep=comma,x=tick,y=mean_abs_excess,restrict expr to domain={\thisrow{series_id}}{30:30}] {\datadir/trajectories.csv};
\addplot[thick,dashed,green!50!black] table[col sep=comma,x=tick,y=mean_abs_excess,restrict expr to domain={\thisrow{series_id}}{40:40}] {\datadir/trajectories.csv};
\legend{E0,E1-baja,E1-alta,E3}
\end{axis}
\end{tikzpicture}
\caption{Exceso de error medio $|d_t|$ por tick, promediado sobre todas las celdas y semillas.}
\label{fig:trayectorias}
\end{figure}

\section{H1: corrección simulada frente al referente analítico}
\label{sec:res-h1}

El cálculo ideal predice 20 ticks hasta la banda con capacidad baja $c=0.005$, todo el horizonte de simulación; la regla de capacidad alta $c=0.03$ requiere solo cuatro ticks antes de que actúen los demás agentes. En las corridas reservadas, el área de exceso media es $2.363$ en E0, $1.056$ en E1-baja y $0.100$ en E1-alta. La diferencia pareada E1-alta menos E0 es $-2.263$ (mediana $-2.078$; intervalo bootstrap por semilla del 95\,\% $[-2.345,-2.182]$), y los promedios pareados de las 300 semillas son menores con E1-alta. La diferencia E1-alta menos E1-baja es $-0.956$ (intervalo $[-1.003,-0.909]$). Estos valores describen cómo el simulador especificado se aparta del referente simple; la dirección de la corrección por arbitraje se deriva de la propia regla.

\begin{table}[htbp]
\centering
\caption{Resúmenes Monte Carlo pareados (300 semillas). Los intervalos describen variación por semilla bajo el simulador fijo.}
\label{tab:inferencia}
\scriptsize
\setlength{\tabcolsep}{3pt}
\begin{tabularx}{\textwidth}{@{}>{\raggedright\arraybackslash}Xrr>{\raggedright\arraybackslash}p{0.25\textwidth}>{\raggedright\arraybackslash}p{0.16\textwidth}@{}}
\toprule
Contraste & Media $\Delta$ & Mediana & Intervalo 95\,\% & Signo \\
\midrule
H1: área E1-alta menos E0 & $-2.2629$ & $-2.0784$ & $[-2.3448,-2.1820]$ & 100\,\% menor \\
H1: área E1-alta menos E1-baja & $-0.9560$ & $-0.8456$ & $[-1.0034,-0.9091]$ & 100\,\% menor \\
H2: desplazamiento máximo legado & $+0.01143$ & $+0.01140$ & $[0.01074,0.01212]$ & 99.67\,\% mayor \\
H3 mapa menos persistencia, MAE & $-0.000808$ & $-0.000960$ & $[-0.000917,-0.000692]$ & 84.67\,\% menor \\
H3 mapa menos AR(1) escalar, MAE & $+0.0000822$ & $+0.0000844$ & $[0.0000780,0.0000864]$ & 97.67\,\% mayor \\
\bottomrule
\end{tabularx}
\end{table}

\subsection{Tasa de recuperación}

La recuperación sostenida depende de la capacidad y de la celda (tabla~\ref{tab:recuperacion}). Sin arbitraje apenas ocurre. Con capacidad baja varía entre 20 y 100\,\%, y es más baja cuanto mayor es el choque: un choque grande exige más pasos acotados por $c\,x$. Con capacidad alta y máxima es total en todas las celdas.

\begin{table}[htbp]
\centering
\caption{Tasa de recuperación sostenida (\%) por celda y condición.}
\label{tab:recuperacion}
\small
\begin{tabular}{@{}rrrrrrr@{}}
\toprule
$\kappa$ & Choque & E0 & E1-baja & E1-alta & E3 & E1-máx \\
\midrule
0.5 & $-0.05$ & 21.7 & 54.7 & 100.0 & 87.0 & 100.0 \\
0.5 & $-0.10$ & 1.3 & 45.3 & 100.0 & 86.7 & 100.0 \\
0.5 & $-0.20$ & 5.3 & 29.3 & 100.0 & 95.7 & 100.0 \\
1 & $-0.05$ & 0.0 & 72.0 & 100.0 & 100.0 & 100.0 \\
1 & $-0.10$ & 3.7 & 59.3 & 100.0 & 100.0 & 100.0 \\
1 & $-0.20$ & 0.0 & 20.3 & 100.0 & 100.0 & 100.0 \\
2 & $-0.05$ & 9.7 & 100.0 & 100.0 & 100.0 & 100.0 \\
2 & $-0.10$ & 0.0 & 99.0 & 100.0 & 100.0 & 100.0 \\
2 & $-0.20$ & 0.0 & 31.0 & 100.0 & 100.0 & 100.0 \\
\bottomrule
\end{tabular}
\end{table}

La tabla~\ref{tab:h1-ideal-times} compara las tasas de retorno simulado al plazo predicho por el referente de piscina aislada con las tasas al horizonte completo de 20 ticks. La primera tasa no representa el evento de primera entrada del referente: la recuperación simulada exige dos observaciones posarbitraje consecutivas dentro de la banda y registra el tick inicial de ese par. La tabla alinea el plazo y conserva el criterio más exigente de retorno sostenido. Cada celda contiene 300 semillas; el CSV reproducible incluye intervalos de Wilson.

\begin{table}[htbp]
\centering
\caption{Tasas de retorno sostenido simulado al plazo ideal de primera entrada y al tick 20. Cada valor indica porcentaje por $N_\gamma$ / porcentaje al tick 20 ($n=300$).}
\label{tab:h1-ideal-times}
\scriptsize
\setlength{\tabcolsep}{4pt}
\begin{tabularx}{\textwidth}{@{}rrr>{\centering\arraybackslash}X>{\centering\arraybackslash}X@{}}
\toprule
Profundidad $\kappa$ & Choque & $N_\gamma$ ($c=.005/.03$) & $c=.005$ & $c=.03$ \\
\midrule
0.5 & $-0.05$ & 5 / 1 & 17.7 / 54.7 & 49.0 / 100.0 \\
1.0 & $-0.05$ & 5 / 1 & 21.3 / 72.0 & 57.3 / 100.0 \\
2.0 & $-0.05$ & 5 / 1 & 43.7 / 100.0 & 100.0 / 100.0 \\
0.5 & $-0.10$ & 10 / 2 & 28.3 / 45.3 & 59.3 / 100.0 \\
1.0 & $-0.10$ & 10 / 2 & 28.7 / 59.3 & 83.3 / 100.0 \\
2.0 & $-0.10$ & 10 / 2 & 44.7 / 99.0 & 100.0 / 100.0 \\
0.5 & $-0.20$ & 20 / 4 & 29.3 / 29.3 & 84.0 / 100.0 \\
1.0 & $-0.20$ & 20 / 4 & 20.3 / 20.3 & 79.7 / 100.0 \\
2.0 & $-0.20$ & 20 / 4 & 31.0 / 31.0 & 98.3 / 100.0 \\
\bottomrule
\end{tabularx}
\end{table}

La tabla también muestra cuándo el referente aislado es informativo y cuándo importan las interacciones. Con capacidad $c=0.03$ y factor de profundidad $\kappa=2$, el 100\,\% de las trayectorias satisface el criterio de retorno sostenido antes del plazo ideal para los choques $-0.05$ y $-0.10$, y el 98.3\,\% lo hace para $-0.20$. Con $\kappa=0.5$ y $1$, solo el 49.0--84.0\,\% y el 57.3--83.3\,\%, respectivamente, lo logra antes de ese plazo, aunque todas las trayectorias con $c=0.03$ se recuperan antes del tick 20. La diferencia refleja las órdenes continuas y el orden de eventos del calendario simulado; no contradice el límite de piscina aislada, cuyos supuestos excluyen esos mecanismos. La simulación cumple así una función acotada: cuantificar las desviaciones en horizonte finito frente a un referente mecánico.

\subsection{Desglose por celda}

El efecto de H1 es negativo en todas las celdas (tabla~\ref{tab:celdas}); su magnitud crece con el choque, porque sin arbitraje el exceso se acumula durante los 20 ticks. La heterogeneidad entre semillas es grande (figura~\ref{fig:cuantiles}): el contraste va de $-4.60$ a $-1.38$.

\begin{table}[htbp]
\centering
\caption{Diferencia media por celda de los contrastes H1 (primario y de capacidad) y H2.}
\label{tab:celdas}
\small
\begin{tabular}{@{}rrrrrc@{}}
\toprule
$\kappa$ & Choque & E1-alta vs. E0 & E1-alta vs. E1-baja & Diferencia en $M$ legado & Reservada \\
\midrule
0.5 & $-0.05$ & $-1.630$ & $-0.832$ & $+0.0523$ & sí \\
0.5 & $-0.10$ & $-1.873$ & $-0.941$ & $+0.0361$ & sí \\
0.5 & $-0.20$ & $-3.128$ & $-1.703$ & $+0.0146$ & sí \\
1 & $-0.05$ & $-1.222$ & $-0.294$ & $0.0000$ & sí \\
1 & $-0.10$ & $-2.023$ & $-0.652$ & $0.0000$ & no \\
1 & $-0.20$ & $-3.784$ & $-1.883$ & $0.0000$ & sí \\
2 & $-0.05$ & $-1.007$ & $-0.116$ & $0.0000$ & sí \\
2 & $-0.10$ & $-1.951$ & $-0.450$ & $0.0000$ & sí \\
2 & $-0.20$ & $-3.748$ & $-1.732$ & $0.0000$ & sí \\
\bottomrule
\end{tabular}
\end{table}

\begin{figure}[htbp]
\centering
\begin{tikzpicture}
\begin{axis}[width=0.8\textwidth,height=6cm,xlabel={Cuantil},ylabel={Contraste por semilla},grid=major,legend pos=south east,legend style={font=\scriptsize}]
\addplot[thick,blue!70!black,mark=*,mark size=1.5pt] table[col sep=comma,x=quantile,y=value,restrict expr to domain={\thisrow{contrast_id}}{1:1}] {\datadir/delta_quantiles.csv};
\addplot[thick,orange!85!black,mark=square*,mark size=1.5pt] table[col sep=comma,x=quantile,y=value,restrict expr to domain={\thisrow{contrast_id}}{2:2}] {\datadir/delta_quantiles.csv};
\legend{H1 primario,H1 capacidad}
\end{axis}
\end{tikzpicture}
\caption{Cuantiles del contraste por semilla de H1 ($\Delta_1$ y comparación de capacidades).}
\label{fig:cuantiles}
\end{figure}

\section{H2: sensibilidad al retiro y a la escala de las órdenes}
\label{sec:res-h2}

Con el estadístico legado $M=\max_t|\log(p_t^D/p_0^D)|$, la diferencia pareada E3 menos E1-alta es $+0.0114$ en promedio (intervalo bootstrap por semilla del 95\,\% $[0.0107,0.0121]$). Este estadístico mide el desplazamiento respecto del precio del AMM previo al choque, no el sobrepaso ni el error máximo de seguimiento respecto de la nueva referencia CEX. Además, se satura en $|\text{choque}|+\gamma$ en seis de nueve celdas. Se conserva solo como descripción del diseño original. El máximo de horizonte completo $P=\max_{0\le t\le20}|e_t|$ tiene una limitación relacionada: incluye el tick 1, cuando el choque impuesto domina ambas condiciones, por lo que su contraste disminuye al crecer el choque.

Las diferencias por celda de la tabla~\ref{tab:celdas} son positivas solo con factor de profundidad $\kappa=0.5$ y nulas con $\kappa=1$ y $2$. El retiro se activa en todas las trayectorias E3 porque cada choque aplicado supera el umbral fijo $|e_t|>0.02$; normalmente ocurre en el tick 1. Esto es una intervención determinista única bajo los choques ensayados, no una prueba de respuestas heterogéneas o graduales. La orden de pánico se activa solo con $\kappa=0.5$, donde el impacto configurado de la orden de ruido es 1.6\,\%, apenas 0.1 puntos porcentuales por encima del umbral de 1.5\,\%.

Por ello, repetimos las trayectorias con semillas archivadas y usamos como diagnóstico exploratorio H2 principal el exceso de error acumulado después de confirmar la recuperación. Sea $t_c$ el segundo de dos ticks consecutivos dentro de la banda que definen el retorno sostenido; entonces $A_{\mathrm{post}}=\sum_{t>t_c}(|e_t|-\gamma)_+$. Con profundidad $\kappa=0.5$, la diferencia pareada E3 menos E1-alta es $0.1229$ (intervalo bootstrap por semilla del 95\,\% $[0.1092,0.1373]$, $n=267$), $0.1239$ ($[0.1134,0.1348]$, $n=272$) y $0.1215$ ($[0.1109,0.1326]$, $n=287$) para choques de $-0.05$, $-0.10$ y $-0.20$, respectivamente. Con profundidades 1 y 2, los contrastes pareados son cero a precisión numérica en todas las celdas de choque. La figura~\ref{fig:h2-post-recovery} muestra este patrón por profundidad. La similitud de los contrastes de baja profundidad entre los tamaños de choque, tras excluir el choque y los dos ticks de confirmación, es más pertinente para estudiar la fragilidad posterior al retorno que la métrica pico anterior. Sigue siendo una comparación posterior al análisis inicial y condicionada a trayectorias recuperadas, no un efecto confirmatorio ni de mercado.

\begin{figure}[htbp]
\centering
\begin{tikzpicture}
\begin{axis}[
  width=0.88\textwidth,height=6.0cm,
  xlabel={Factor de profundidad $\kappa$}, ylabel={$\Delta A_{\mathrm{post}}$ (E3 $-$ E1-alta)},
  xmin=0.25,xmax=2.25,ymin=-0.01,ymax=0.16,
  xtick={0.5,1,2}, ymajorgrids=true, grid style={gray!20},
  legend style={font=\scriptsize,at={(0.98,0.98)},anchor=north east},
  axis line style={gray!60}, tick label style={font=\scriptsize}
]
\addplot+[only marks,mark=*,mark size=2pt,blue!70!black,error bars/.cd,y dir=both,y explicit] coordinates {
  (0.42,0.12295) +- (0,0.01402) (0.92,0) (1.92,0)
};
\addlegendentry{Choque $-0.05$}
\addplot+[only marks,mark=square*,mark size=2pt,orange!85!black,error bars/.cd,y dir=both,y explicit] coordinates {
  (0.50,0.12393) +- (0,0.01071) (1.00,0) (2.00,0)
};
\addlegendentry{Choque $-0.10$}
\addplot+[only marks,mark=triangle*,mark size=2.2pt,green!45!black,error bars/.cd,y dir=both,y explicit] coordinates {
  (0.58,0.12154) +- (0,0.01083) (1.08,0) (2.08,0)
};
\addlegendentry{Choque $-0.20$}
\end{axis}
\end{tikzpicture}
\caption{Contraste del área de exceso posterior a la recuperación según la profundidad probada. Las barras aproximan de forma simétrica los intervalos bootstrap pareados del 95\,\% con $\kappa=0.5$; los contrastes con $\kappa=1$ y $2$ son cero a precisión numérica. En baja profundidad, los tamaños pareados son 267, 272 y 287; las tres profundidades usan 300 semillas archivadas.}
\label{fig:h2-post-recovery}
\end{figure}

El contraejemplo de la razón de presión se puede interpretar mecánicamente dentro del diseño ensayado. Con profundidad baja y capacidad $c=0.03$, el flujo continuo configurado de ruido más tendencia produce $R_{\mathrm{flow}}=0.670$ en el caso estático, por debajo de uno; la orden de pánico única produce $R_{\mathrm{panic}}=1.327$, y el flujo adverso máximo configurado de ruido más pánico da $R_{\max}^{\mathrm{adv}}=1.586$. En la reserva nominal de E3 tras el retiro, las razones son 1.208, 2.375 y 2.831 (tabla~\ref{tab:pressure-ratio}). La repetición sin pánico no muestra salidas posteriores en E1-alta estático (0/300 para cada choque; límite superior Wilson del 95\,\% de 1.3\,\%), mientras que las salidas con pánico permitido son episodios únicos y las 143 coinciden con el tick de la orden de pánico. En E3, el 98.2--99.6\,\% de las trayectorias recuperadas vuelve a salir, y la mayoría presenta episodios múltiples. Este patrón se mantiene al desactivar el pánico. En estas celdas, la evidencia respalda una salida única asociada al pánico para E1 y salidas recurrentes bajo E3 que no requieren pánico. No establece umbrales universales: las razones usan órdenes configuradas y reservas aproximadas, no flujo direccional realizado en cada evento.

El control pareado de referencia de capacidad con $\kappa=0.5$ fija el límite de entrada del arbitraje en $c$ por la reserva inicial, mientras la ejecución sigue limitada por las reservas activas. Sus razones con límite fijo coinciden con las de E1 estático ($R_{\mathrm{flow}}\approx0.67$, $R_{\mathrm{panic}}\approx1.32$), y el contraste del área posrecuperación disminuye de 0.122--0.124 a 0.0036--0.0097, como predice $R$. Esto no descompone efectos independientes de capacidad y profundidad: ambas son componentes de la misma razón orden/reserva. Como las E3 que no se recuperan quedan fuera del contraste condicional, la estimación puede subestimar la fragilidad sobre todas las trayectorias. El diseño factorial más amplio que varíe retiro, capacidad y profundidad sigue siendo necesario.



El seguimiento amplio requiere un diseño factorial prerregistrado que varíe por separado retiro, referencia de capacidad, profundidad, escala de órdenes y pánico. El replay actual es exploratorio y reutiliza semillas existentes; no aporta una muestra independiente.

\section{H3: comparación predictiva con referencias simples}
\label{sec:res-h3}

La cláusula del mapa se responde aquí comparando la descripción predictiva del error continuo en trayectorias reservadas del simulador; las etiquetas de horizonte finito son resúmenes descriptivos, no clases observadas. El mapa de dos estados se ajusta al exceso de error con signo $d_t$, que vale cero dentro de la banda y cambia de pendiente en su límite. Una sola relación afín no puede representar exactamente este umbral. En celdas reservadas, el mapa tiene menor error absoluto medio que la persistencia en la mayoría de condiciones, pero no supera al AR(1) escalar en ninguna condición principal. Por tanto, su segundo estado no demuestra valor predictivo incremental a un paso en este diseño; el mapa queda como resumen descriptivo del simulador, no como modelo de regímenes.

\begin{table}[htbp]
\centering
\caption{Error absoluto medio a un paso en celdas reservadas del simulador.}
\label{tab:h3}
\small
\begin{tabular}{@{}lrrrc@{}}
\toprule
Condición & Mapa dos estados & Persistencia & AR(1) escalar & Mapa mejora persistencia \\
\midrule
E0 & 0.01765 & 0.01751 & 0.01757 & no \\
E1-baja & 0.01573 & 0.01615 & 0.01566 & sí \\
E1-alta & 0.00493 & 0.00620 & 0.00486 & sí \\
E3 & 0.00929 & 0.01096 & 0.00916 & sí \\
E1-máx$^{\ast}$ & 0.00567 & 0.00594 & 0.00578 & sí \\
\bottomrule
\multicolumn{5}{@{}l}{\footnotesize $^{\ast}$ Condición de extrapolación fuera de los niveles principales de capacidad.}
\end{tabular}
\end{table}

Promediada sobre las cuatro condiciones principales, la diferencia mapa menos persistencia es $-0.000808$ (intervalo bootstrap por semilla del 95\,\% $[-0.000917,-0.000692]$); la diferencia mapa menos AR escalar es $+0.0000822$ (intervalo $[0.0000780,0.0000864]$). Estos intervalos describen variación Monte Carlo. Las etiquetas de trayectorias finitas resumen las mismas corridas y no son regímenes observados de forma independiente.

\section{Comparación reservada de mapas afines e híbridos}
\label{sec:hybrid-results}

La nueva comparación usa el análisis 0.1.0, ejecución \code{\HybridRunID} y el error de precio con signo. En 9\,600 trayectorias reservadas de 300 semillas nuevas, el mapa híbrido redujo el MAE recursivo medio del error de precio de 0.030807 a 0.030291. La mejora pareada media fue 0.000516 (intervalo bootstrap por semilla del 95\,\% [0.000499, 0.000533]); el 99.7\,\% de los contrastes promediados por semilla favoreció a Hyb. La mejora media frente a L0 es de aproximadamente 1.7\,\%. Es una pequeña ganancia en error de reproducción dentro del simulador especificado; la evaluación recibe la trayectoria futura realizada del CEX y no constituye un pronóstico en tiempo real.

Cada condición aportó 12\,000 transiciones de selección. Hubo observaciones a ambos lados de la banda (tabla~\ref{tab:hybrid-regimes}), aunque el modelo de arbitraje bajo tuvo menos observaciones dentro de ella. El ajuste sigue siendo específico para cada condición, y estos conteos no demuestran que las dos variables de estado capturen toda la memoria relevante para la transición.

\begin{table}[htbp]
\centering
\caption{Transiciones de selección dentro y fuera de la banda configurada.}
\label{tab:hybrid-regimes}
\small
\begin{tabular}{@{}lrrr@{}}
\toprule
Condición & En banda & Fuera de banda & Total \\
\midrule
E0 & 681 & 11\,319 & 12\,000 \\
E1-baja & 1\,466 & 10\,534 & 12\,000 \\
E1-alta & 3\,776 & 8\,224 & 12\,000 \\
E3 & 2\,834 & 9\,166 & 12\,000 \\
\bottomrule
\end{tabular}
\end{table}

La mejora no se distribuye entre las condiciones operativas. Se concentra en el arbitraje de capacidad baja, especialmente con profundidad 2, donde la ganancia promediada entre los cuatro choques reservados es 0.003943. Con profundidad 0.5, la ganancia correspondiente es 0.000354. Los promedios de las demás combinaciones de condición y profundidad se encuentran entre aproximadamente $-0.000101$ y $-0.000002$, lo que favorece levemente a L0. El MAE secundario a un paso para $e$ mejora de 0.014882 a 0.014800, mientras que los errores a un paso y recursivos de la presión de flujo $z$ son ligeramente peores con Hyb. Por tanto, el resultado respalda una ventaja limitada y dependiente de la condición de los términos de banda para predecir el error de precio; no muestra una superioridad uniforme del mapa híbrido ni que las variables de estado representen todo el proceso.

\begin{figure}[H]
\centering
\begin{tikzpicture}
\begin{axis}[
  xbar, width=0.93\textwidth, height=6.1cm,
  xmin=-0.0002, xmax=0.0045,
  xlabel={Ganancia del MAE recursivo reservado, L0 menos Hyb},
  symbolic y coords={C1,C2,C3,C4,C5,C6,C7,C8}, ytick=data,
  yticklabels={E0, E1-baja, E1-alta, E3, E0, E1-baja, E1-alta, E3},
  yticklabel style={font=\scriptsize},
  y dir=reverse, xmajorgrids=true, grid style={gray!20},
  axis x line*=bottom, axis y line*=left,
  title style={font=\small}, title={Profundidad 0.5 (arriba) y profundidad 2 (abajo); valores positivos favorecen Hyb}
]
\addplot[fill=blue!55!black,draw=blue!55!black] coordinates {
  (-0.000002,C1) (0.000354,C2) (-0.000014,C3) (-0.000003,C4)
  (-0.000017,C5) (0.003943,C6) (-0.000033,C7) (-0.000101,C8)
};
\end{axis}
\end{tikzpicture}
\caption{Ganancia media del MAE recursivo del error de precio por condición operativa y profundidad. Cada barra promedia cuatro magnitudes de choque reservadas; los contrastes por celda son descriptivos y no tienen intervalos de confianza separados. Los valores se exportan en \code{figures/data/hybrid\_depth\_condition.csv}.}
\label{fig:hybrid-cell-gains}
\end{figure}

La ganancia media del mapa híbrido no debe interpretarse aisladamente: casi toda la ganancia proviene de una combinación de baja capacidad y profundidad, mientras que la predicción de la presión de flujo no mejora. El siguiente paso de modelado consiste en representar explícitamente los límites de capacidad y las transiciones de retiro de LP y evaluarlos con otro protocolo congelado, en vez de añadir más términos polinómicos libres. La comparación de mapas es un diagnóstico de reproducción dentro del simulador, no un pronóstico de mercado ni una afirmación de estabilidad del sistema completo.

\section{Trayectorias por profundidad}

La figura~\ref{fig:prof} compara E1-alta y E3 con $\kappa=0.5$. Con la menor profundidad, E3 parte del mismo exceso inicial y se corrige más despacio. Una explicación compatible es que el retiro reduce tanto la profundidad como el límite de arbitraje $c\,x$; estos mecanismos están acoplados en el diseño actual. Con $\kappa=1$ y $2$ las dos trayectorias son casi idénticas. El retorno sostenido ocurre en promedio en el tick $7.8$ para E3 con $\kappa=0.5$ frente a $2.8$ para E1-alta, mientras que con $\kappa=1$ y $2$ las diferencias son de $0.1$ ticks o menos.

\begin{figure}[htbp]
\centering
\begin{tikzpicture}
\begin{axis}[width=0.9\textwidth,height=6.2cm,xlabel={Tick},ylabel={$|d_t|$ medio},grid=major,legend style={font=\scriptsize},xmin=0,xmax=20,ymin=0]
\addplot[thick,blue!70!black] table[col sep=comma,x=tick,y=mean_abs_excess,restrict expr to domain={\thisrow{series_id}}{31:31}] {\datadir/trajectories.csv};
\addplot[thick,dashed,green!50!black] table[col sep=comma,x=tick,y=mean_abs_excess,restrict expr to domain={\thisrow{series_id}}{41:41}] {\datadir/trajectories.csv};
\legend{E1-alta,E3}
\end{axis}
\end{tikzpicture}
\caption{Exceso de error medio con $\kappa=0.5$: arbitraje alto estático (E1-alta) y con retiro por umbral (E3).}
\label{fig:prof}
\end{figure}

\section{Sensibilidad de ingeniería: recuperación y escala de las órdenes}
\label{sec:engineering-sensitivity}

Un análisis suplementario (análisis 0.6.0, ejecución \code{\EngineeringRunID}) usó 300 semillas nuevas (2000--2299), tres variantes de flujo de órdenes y las 15 celdas de profundidad/choque/condición de la grilla de diseño reservada: 40\,500 trayectorias de 20 ticks abstractos. Es un análisis descriptivo y exploratorio de sensibilidad, no un contraste inferencial. La recuperación significa el primer par observado de ticks consecutivos dentro de la banda; las salidas posteriores se contabilizan por separado. Los plazos se expresan en ticks, no en tiempo transcurrido.

Con los tamaños absolutos originales, E3 con retiro por umbral se recuperó al menos una vez en el 97.4\,\% de las trayectorias, pero solo el 81.9\,\% terminó dentro de la banda. De las trayectorias que se recuperaron, el 31.2\,\% salió después. Con profundidad 0.5, solo el 45.8\,\% terminó dentro de la banda y el 98.8\,\% de las trayectorias recuperadas salió posteriormente. Por tanto, en este modelo, un primer retorno sostenido no equivale a un seguimiento persistente. Estos promedios igualmente ponderados del diseño no sustentan una afirmación poblacional con ajuste por multiplicidad.

El contraste de respuesta máxima entre E3 con retiro por umbral y estático cambia según la definición de las órdenes (tabla~\ref{tab:engineering-sensitivity}). Mantener fijos los tamaños absolutos produjo un efecto con profundidad 0.5; al escalar las órdenes de ruido y tendencia con las reservas iniciales, el contraste fue efectivamente cero en las tres profundidades. Desactivar el pánico redujo, pero no eliminó, esta métrica de horizonte completo dominada por el choque. El análisis posterior a la recuperación que sigue separa ese evento único de las salidas posteriores. Esta sensibilidad muestra que la profundidad por sí sola no determina el resultado; también importan el tamaño de la orden relativo a las reservas y la regla de reacción.

\begin{table}[htbp]
\centering
\caption{Respuesta máxima suplementaria pareada: E3 con retiro por umbral menos estático por profundidad (300 semillas nuevas). Intervalos aproximados del 95\,\% sin ajuste.}
\label{tab:engineering-sensitivity}
\small
\begin{tabular}{@{}lrrr@{}}
\toprule
Variante del flujo de órdenes & Profundidad 0.5 & Profundidad 1 & Profundidad 2 \\
\midrule
Órdenes absolutas fijas & 0.03458 [0.03225, 0.03690] & 0.00000 & 0.00000 \\
Órdenes proporcionales a reservas iniciales & 0.00000 & 0.00000 & 0.00000 \\
Órdenes fijas, sin pánico & 0.01471 [0.01364, 0.01577] & 0.00000 & 0.00000 \\
\bottomrule
\end{tabular}
\end{table}

Con profundidad 0.5 y órdenes absolutas fijas, E3 con retiro por umbral terminó dentro de la banda en el 45.8\,\% de las trayectorias; con órdenes proporcionales, terminó dentro en el 100\,\%. Estas variantes comparten el mismo modelo estilizado de agentes representativos y no constituyen comparaciones empíricas. Los resúmenes por celda, contrastes pareados, configuración y hashes completos se exportan en \code{figures/engineering/}.

\subsection{Presión y mecanismos de salida: repetición exploratoria}

El protocolo exploratorio congelado 0.7.0 repite 9\,000 trayectorias con semillas existentes 2000--2299: 5\,400 con órdenes fijas en las tres profundidades y choques, 1\,800 trayectorias poco profundas sin pánico y 1\,800 trayectorias poco profundas con el límite de entrada del arbitraje fijado a su valor basado en reservas iniciales. Las 7\,200 trayectorias con el límite ligado a las reservas originales reproducen exactamente las métricas archivadas de recuperación, salidas, área de exceso, respuesta máxima y límites de capital. Es un análisis de reproducibilidad y mecanismos con semillas existentes, no una muestra Monte Carlo independiente.

Para una orden configurada $q$ y reserva $x$, definimos $R(q,x,c)=2\log(1+q/x)/\delta_c$, donde $\delta_c=2\log(1+c)$. El diagnóstico conservador de un tick es $R_{\mathrm{1tick}}=\max_{q\in\mathcal Q_{\mathrm{adv}}}R(q,x,c)$, donde $\mathcal Q_{\mathrm{adv}}$ incluye la mayor entrada adversa configurada, incluido el pánico. Con profundidad 0.5 y capacidad $c=0.03$, el desglose informa $R_{\mathrm{flow}}$ para $q=100$ (ruido y tendencia alineados), $R_{\mathrm{panic}}$ para la orden única de pánico $q=200$, y $R_{\mathrm{1tick}}=R_{\max}^{\mathrm{adv}}$ para $q=240$ (ruido y pánico alineados; el pánico sustituye a la tendencia, no se suma a ella). Para E3, la reserva posterior al retiro se aproxima como el 55\,\% de la inicial. Son razones configuradas, no mediciones de flujo por tick: el ruido ocurre con probabilidad 0.7, las direcciones pueden diferir y las reservas cambian durante el tick.

\begin{table}[htbp]
\centering
\caption{Razones configuradas de flujo a capacidad con factor de profundidad 0.5 y $c=0.03$.}
\label{tab:pressure-ratio}
\small
\begin{tabular}{@{}lrrr@{}}
\toprule
Condición & $R_{\mathrm{flow}}$ ($q=100$) & $R_{\mathrm{panic}}$ ($q=200$) & $R_{\max}^{\mathrm{adv}}$ ($q=240$) \\
\midrule
E1-alta, estático ($x\approx5000$) & 0.670 & 1.327 & 1.586 \\
E3, tras retiro del 45\,\% ($x\approx2750$) & 1.208 & 2.375 & 2.831 \\
E3, límite fijo en reservas iniciales ($x\approx2750$) & 0.673 & 1.322 & 1.575 \\
\bottomrule
\end{tabular}
\end{table}

La tabla~\ref{tab:departure-episodes} separa si una trayectoria sale después de la recuperación de cuántas veces lo hace. En E1-alta estático, las 90 salidas después del choque $-0.05$ y las 53 después de $-0.10$ son episodios únicos, y cada primera salida coincide con el tick de la orden de pánico; no hay salidas tras la ablación pareada sin pánico para ninguno de los choques (0/300 en cada celda; límite superior Wilson del 95\,\% de 1.3\,\%). En E3, el 98.2--99.6\,\% de las trayectorias recuperadas vuelve a salir posteriormente, y la mayoría tiene episodios múltiples. Este patrón se mantiene en las corridas sin pánico. En estas celdas, la evidencia respalda un mecanismo E1 de salida única asociada al pánico y un mecanismo E3 de flujo recurrente que no requiere pánico; no convierte las razones aproximadas en umbrales universales.

\begin{table}[htbp]
\centering
\caption{Trayectorias que salen de la banda y episodios tras la recuperación con profundidad 0.5. El valor entre paréntesis es el promedio de episodios consecutivos fuera de banda por trayectoria recuperada.}
\label{tab:departure-episodes}
\scriptsize
\setlength{\tabcolsep}{4pt}
\begin{tabularx}{\textwidth}{@{}l r >{\centering\arraybackslash}X >{\centering\arraybackslash}X@{}}
\toprule
Condición & Choque & Pánico activo: trayectorias/recuperadas (episodios/trayectoria) & Pánico desactivado: trayectorias/recuperadas (episodios/trayectoria) \\
\midrule
E1-alta estático & $-0.05$ & 90/300 (0.30) & 0/300 (0.00) \\
E1-alta estático & $-0.10$ & 53/300 (0.18) & 0/300 (0.00) \\
E1-alta estático & $-0.20$ & 0/300 (0.00) & 0/300 (0.00) \\
E3 retiro por umbral & $-0.05$ & 266/267 (4.09) & 279/280 (4.93) \\
E3 retiro por umbral & $-0.10$ & 271/272 (4.41) & 256/258 (4.09) \\
E3 retiro por umbral & $-0.20$ & 284/287 (3.65) & 288/288 (3.77) \\
\bottomrule
\end{tabularx}
\end{table}

El control pareado de referencia de capacidad fija el límite de entrada de las órdenes de arbitraje en el 3\,\% de la reserva inicial del activo de entrada y conserva el límite de ejecución por reservas activas. El control iguala las razones configuradas con el caso E1 estático ($R_{\mathrm{flow}}\approx0.67$ y $R_{\mathrm{panic}}\approx1.32$); el contraste del área posrecuperación cae de 0.122--0.124 a 0.0036--0.0097, como predice el diagnóstico de razones. Esto no descompone el efecto en contribuciones independientes de «capacidad» y «profundidad»: ambas entran en la misma razón orden/reserva y quedan acopladas por construcción. Como el contraste se condiciona a trayectorias que se recuperan, excluir las E3 que no se recuperan probablemente sesga esta estimación condicional hacia abajo frente a una estimación de fragilidad sobre todas las trayectorias. El diseño factorial más amplio que varíe retiro, capacidad y profundidad sigue siendo necesario.

\begin{table}[htbp]
\centering
\caption{Control de referencia de capacidad en la celda poco profunda. La recuperación es la proporción de las 300 trayectorias. La salida se informa como salidas/trayectorias recuperadas con seguimiento posterior a la confirmación, igual que en la tabla~\ref{tab:departure-episodes}; $\Delta A_{\mathrm{post}}$ es E3 menos E1-alta entre las mismas trayectorias pareadas recuperadas.}
\label{tab:capacity-reference}
\scriptsize
\setlength{\tabcolsep}{4pt}
\begin{tabularx}{\textwidth}{@{}r >{\centering\arraybackslash}X >{\centering\arraybackslash}X rr@{}}
\toprule
Choque & Recuperación E3: límite ligado/fijo & Salida E3: límite ligado/fijo & $\Delta A_{\mathrm{post}}$ ligado & fijo \\
\midrule
$-0.05$ & 90.3/100.0\,\% & 266/267\,(99.6)/123/300\,(41.0) & 0.1229 & 0.0097 \\
$-0.10$ & 91.0/100.0\,\% & 271/272\,(99.6)/101/300\,(33.7) & 0.1239 & 0.0080 \\
$-0.20$ & 95.7/100.0\,\% & 284/287\,(99.0)/53/300\,(17.7) & 0.1215 & 0.0036 \\
\bottomrule
\end{tabularx}
\end{table}

En la variante de capacidad fija con pánico activo, las 123, 101 y 53 trayectorias que salen contienen exactamente un episodio cada una, y los 277 episodios coinciden con el tick del pánico. La comprobación dirigida de 900 trayectorias con capacidad fija y sin pánico no produjo ninguna salida posterior en ninguna celda de choque (0/900 por celda); todas se recuperaron con seguimiento posterior a la confirmación. Estos resultados respaldan $R$ como hipótesis organizadora y preregistrable para el próximo factorial, no como una ley validada.

El seguimiento amplio exige un diseño factorial prerregistrado que varíe por separado retiro, referencia de capacidad, profundidad, escala de órdenes y pánico. La repetición actual es exploratoria y reutiliza semillas; no aporta una muestra independiente.

\section{Resumen de los resultados}

\begin{itemize}
  \item \textbf{Referente analítico del arbitraje:} una piscina aislada con capacidad $c$ reduce el error logarítmico fuera de la banda hasta en $2\log(1+c)$ por tick. La tabla con plazos equiparados muestra lo poco frecuente que es iniciar un retorno sostenido de dos ticks antes del plazo ideal de primera entrada con capacidad baja; la dirección de la corrección sigue incorporada en la regla.
  \item \textbf{Sensibilidad al retiro:} el control con capacidad inicial fija iguala las razones configuradas $R$ con E1 poco profundo y elimina el contraste posrecuperación como se predice. No separa capacidad y profundidad, que entran en la misma razón orden/reserva; el resultado sigue siendo posterior al análisis inicial y específico de esta condición.
  \item \textbf{Mecanismo de las salidas:} en E1-alta estático, las salidas posteriores son episodios únicos asociados al pánico y desaparecen al desactivar esa orden. E3 presenta episodios repetidos en la mayoría de trayectorias recuperadas, incluso sin pánico. Son resultados del modelo programado, no frecuencias de mercado.
  \item \textbf{Razones de presión:} las trayectorias con capacidad fija reproducen episodios únicos alineados con el pánico, mientras que 900 trayectorias sin pánico no presentan salidas posteriores. Esto respalda $R$ como hipótesis organizadora para preregistro, no como ley de umbral ni descomposición entre capacidad y profundidad.
  \item \textbf{Modelos reducidos:} el mapa de dos estados mejora la persistencia en algunas condiciones, pero no supera al AR(1) escalar. La ganancia de cerca de 1.7\,\% del mapa sensible a la banda se concentra en una condición y utiliza la trayectoria futura realizada del CEX.
\end{itemize}

\chapter{Discusión y conclusiones}
\label{cap:discusion}

\section{Discusión de los resultados}

\subsection{Seguimiento de precios y capacidad de arbitraje (OE2, H1)}

El resultado de H1 es el más sólido del trabajo: el arbitraje activo reduce el área de exceso de error en todas las semillas y en todas las celdas, y la capacidad gradúa la recuperación, de $4.6\,\%$ sin arbitraje a $56.8\,\%$ con capacidad baja y a $100\,\%$ con capacidad alta. Es un resultado esperable desde la mecánica del modelo, porque el arbitrajista es el único agente que mueve el precio del AMM hacia la referencia, y por eso la contribución no está en el signo del efecto sino en su graduación. La tasa de recuperación con capacidad baja depende de la combinación entre el tamaño del choque y la profundidad (tabla~\ref{tab:recuperacion}): un choque grande requiere varios pasos acotados, y el horizonte de 20 ticks puede no bastar. Esa dependencia sugiere que, en este modelo, la capacidad de arbitraje relevante no es un umbral sino una razón entre el tamaño de la discrepancia y la capacidad por tick.

La lección de diseño es la precisión. La semiamplitud objetivo de H1 no se alcanzó porque el piloto combinó choques positivos y negativos cuyos efectos se cancelan casi por completo dentro de cada semilla (correlación $-0.981$). Un piloto construido con el mismo diseño que la confirmación habría anticipado una desviación nueve veces mayor. La conclusión de que el arbitraje reduce el exceso no depende de ello, pero la estimación del tamaño del efecto es menos precisa de lo previsto.

\subsection{Liquidez adaptativa (OE3, H2)}

H2 se sostiene en promedio, pero no de manera uniforme. El retiro se activó en todas las ejecuciones de E3 y, sin embargo, la respuesta máxima solo aumenta con la menor profundidad. Con las profundidades 1 y 2 el máximo es el mismo en E1-alta y en E3. Con la profundidad 0.5, la comprobación exploratoria sin agente de pánico conserva aproximadamente la mitad del efecto, de modo que este depende del retiro y, en igual medida, de la interacción entre el retiro y el agente de pánico. Lo que el experimento permite afirmar es que, en este micromundo, una regla de retiro de un solo disparo no amplifica la respuesta máxima cuando las órdenes son pequeñas respecto de la reserva, y que la amplificación aparece cuando la piscina es suficientemente superficial para que una orden de ruido active la reacción de pánico.

No se afirma que los proveedores reales amplifiquen las perturbaciones. La regla de retiro es una decisión de modelado, sin calibración y sin respaldo empírico propio. Su respaldo teórico es indirecto: la literatura trata la salida de los proveedores como un problema de decisión \cite{bergault2025exit} y documenta su comportamiento en intercambios descentralizados \cite{lehar2025uniswap}, pero ninguno de esos trabajos valida la regla usada aquí.

\subsection{Mapa dinámico reducido (OE4, H3)}

El resultado de H3 es mixto. El mapa de dos variables predice mejor que la persistencia, y con una mejora sustancial en las condiciones con arbitraje activo (15 a 21 por ciento en E1-alta y E3), pero no supera a un autorregresivo escalar y es ligeramente peor que la persistencia sin arbitraje. Dos observaciones limitan su interpretación. Primera, el segundo estado, la presión neta de compra, no mejora la predicción a un paso respecto del exceso de error solo; la ventaja del mapa sobre la persistencia proviene del coeficiente de contracción $a$, que un modelo escalar captura igual. Segunda, el coeficiente $\beta$ de E3 tiene signo opuesto al de las demás condiciones y la ecuación del exceso es degenerada en E1-máx, lo que muestra que la estructura de dos variables no es estable entre condiciones.

La clasificación de regímenes con el criterio de Schur no discrimina: todas las matrices ajustadas son estables, incluida la de E0, cuyo radio espectral es $0.998$ y cuyas ejecuciones son persistentes en $57\,\%$ de los casos. Esto es coherente con lo que el criterio puede decir (la estabilidad de una dinámica lineal autónoma en un punto) y no con el horizonte finito y el ruido del simulador. El mapa sirve como descripción de la tasa de contracción del error, no como clasificador de regímenes.

\section{Conclusiones por objetivo}

\begin{description}
  \item[OE1.] Se formuló y verificó un simulador discreto con nueve invariantes y un registro reproducible: la ejecución confirmatoria se regenera con identificador y sumas de verificación idénticos, salvo la marca de tiempo. La verificación demuestra que la implementación respeta su especificación, no que la especificación sea realista.
  \item[OE2.] El arbitraje activo reduce el área de exceso de error de $2.36$ a $0.10$, y la capacidad gradúa la tasa de recuperación en las condiciones simuladas.
  \item[OE3.] El retiro adaptativo aumenta la respuesta máxima en $0.011$ en promedio, pero solo con la menor profundidad ensayada, y aproximadamente la mitad de ese efecto se atribuye al agente de pánico. No se concluye amplificación con profundidades mayores.
  \item[OE4.] El mapa reducido supera a la persistencia, no supera a un autorregresivo escalar, y su clasificación de regímenes fuera de muestra es débil. Es una descripción local, y su utilidad queda limitada a la tasa de contracción del error.
\end{description}

\section{Conclusión general}

En el micromundo estudiado, la capacidad de arbitraje es el determinante principal del seguimiento del precio de referencia por un AMM de producto constante. La profundidad modifica el impacto de las órdenes de forma analítica y, junto con una regla de reacción de los agentes, puede amplificar la respuesta. Un mapa de dos variables describe la contracción del error mejor que la persistencia, pero un modelo de una variable lo describe igual de bien. Estas conclusiones valen para el modelo y no se extienden a mercados reales.

\section{Limitaciones y amenazas a la validez}

\begin{itemize}
  \item \textbf{Referencia exógena.} El CEX no responde al AMM, ni tiene profundidad, latencia o comisiones. Los efectos de realimentación entre mercados quedan fuera.
  \item \textbf{Sin comisiones.} El simulador solo implementa intercambios sin comisión. La banda de no arbitraje es un parámetro y no una consecuencia de un costo. Los niveles de comisión, de latencia y de volatilidad externa son candidatos para trabajos posteriores.
  \item \textbf{Calendario.} El orden de las acciones dentro de un tick afecta los resultados (por ejemplo, el agente de pánico observa el precio antes de que actúen el proveedor y el arbitrajista). No se evaluaron calendarios alternativos.
  \item \textbf{Definición de profundidad.} La profundidad escala las reservas y mantiene constante el tamaño absoluto de las órdenes. Otra definición, con órdenes proporcionales a la reserva, podría cambiar la conclusión sobre H2.
  \item \textbf{Inventarios.} Los inventarios son amplios y ninguna orden estuvo limitada por capital; las restricciones de capital reales no se representan.
  \item \textbf{Agentes de un solo tipo.} Cada clase de agente es representativa y no hay heterogeneidad dentro de una clase.
  \item \textbf{Solo datos sintéticos.} No hay datos de mercado ni validación externa. La coherencia entre las series simuladas y los hechos estilizados de mercados reales no se evaluó.
  \item \textbf{Comprobación exploratoria.} La repetición de H2 sin agente de pánico no estaba preregistrada.
  \item \textbf{Generalización de H3.} La evaluación fuera de muestra cubre celdas y semillas no vistas, pero dentro del mismo simulador y de la misma familia de choques.
\end{itemize}

\chapter{Trabajos futuros}
\label{cap:futuro}

Las líneas siguientes se derivan de las limitaciones del capítulo~\ref{cap:discusion}.

\begin{enumerate}
  \item \textbf{Niveles de comisión, latencia y volatilidad externa.} El simulador debe incorporar el intercambio con comisión de la ecuación~\eqref{eq:comision}, un retraso en la ejecución del arbitraje y una referencia con volatilidad, para convertir en resultados los factores que hoy son candidatos de diseño (tabla~\ref{tab:factores}). Con comisión, la banda de no arbitraje deja de ser un parámetro libre y queda determinada por $\gamma=-\log(1-f)$.
  \item \textbf{CEX con profundidad y latencia.} Sustituir la referencia exógena por un libro de órdenes que responda a las operaciones del AMM permitiría estudiar la realimentación bidireccional.
  \item \textbf{Liquidez concentrada.} Extender el AMM a posiciones por intervalos de precio, que cambian la relación entre reserva e impacto.
  \item \textbf{Reglas de proveedores y escala de órdenes.} Contrastar con datos de eventos reglas graduales o estimadas de respuesta de los proveedores y tamaños de orden ligados a las reservas; la sensibilidad actual fija reglas estilizadas y no identifica la conducta real de los proveedores.
  \item \textbf{Planificación de precisión.} Usar un piloto que replique el diseño de choques con signo y estratificar por dirección; el piloto original subestimó la varianza de H1.
  \item \textbf{Mapa reducido.} Probar mapas con más variables o con términos no lineales, y evaluar si mejoran al autorregresivo escalar; explorar otros criterios de clasificación de regímenes.
  \item \textbf{Validación de ingeniería con un par real.} Definir de antemano un par y un periodo; alinear la mejor compra/venta y profundidad del CEX con las reservas, intercambios y eventos de liquidez del DEX; registrar comisiones, gas, demora de ejecución, orientación de tokens, orden de bloques y derechos de uso de datos; después probar persistencia de recuperación, tamaño relativo de órdenes y desempeño reservado del modelo sustituto. No transferir directamente umbrales sintéticos a producción.
  \item \textbf{Evaluación de ingeniería y pedagogía.} Medir si el tablero mejora la capacidad de operadores o estudiantes para identificar recuperación transitoria y estrés relativo de órdenes/reservas; mantener esta evaluación separada de la validación de mercado.
\end{enumerate}


\bibliographystyle{unsrtnat}
\bibliography{../references}

\chapter{Anexos técnicos}
\label{cap:anexos}

\section{Descripción ODD del modelo}
\label{sec:odd}

La descripción sigue el protocolo ODD \cite{grimm2020odd}.

\subsection*{Visión general}

\textbf{Propósito.} Estudiar cómo un AMM de producto constante sigue un precio de referencia exógeno después de una perturbación, bajo distintas capacidades de arbitraje, profundidades y reglas de provisión de liquidez.

\textbf{Entidades y variables de estado.} El AMM (reservas $x_t$, $y_t$); la referencia CEX ($\pC_t$); cuatro billeteras: ruido, tendencia y pánico, arbitrajista y proveedor de liquidez, con inventarios en $X$ e $Y$. Las variables derivadas son $e_t$, $d_t$ y $z_t$.

\textbf{Escalas.} Un tick es una unidad discreta de tiempo sin equivalencia calendario; el horizonte es de 20 ticks. Las reservas iniciales son $10\,000$ unidades por activo y el precio inicial es 1.

\textbf{Resumen del proceso.} En cada tick: choque de la referencia (solo en el tick 1), orden de ruido, decisión de tendencia o pánico, regla de retiro por umbral, orden de arbitraje, registro y comprobaciones.

\subsection*{Conceptos de diseño}

\textbf{Emergencia.} La trayectoria del error de seguimiento resulta de la interacción entre el arbitraje acotado y el flujo de órdenes. \textbf{Respuestas basadas en reglas.} El proveedor retira liquidez ante un error mayor que $0.02$ (una vez); el agente de pánico reacciona ante un retorno del AMM menor o igual a $-1.5\,\%$ (una vez por ejecución); el agente de tendencia sigue retornos de magnitud mayor o igual a $0.1\,\%$. \textbf{Objetivos.} Ninguno de los agentes optimiza; las reglas son heurísticas. \textbf{Observación.} Los agentes observan el precio del AMM y, el arbitrajista, también la referencia. \textbf{Estocasticidad.} Solo el flujo de ruido es aleatorio (actividad y lado); la mecánica del AMM y las reglas son deterministas. \textbf{Interacción.} Indirecta, a través del precio y de las reservas.

\subsection*{Detalles}

\textbf{Inicialización.} Reservas $10\,000/10\,000$, precio inicial de la referencia 1; billeteras de ruido y de tendencia con dos veces las reservas iniciales, arbitrajista con cinco veces, proveedor con cero. \textbf{Entrada.} El choque de la referencia y el calendario de ruido pregenerado a partir de la semilla. \textbf{Submodelos.} Las ecuaciones~\eqref{eq:reservas} a~\eqref{eq:qstar} y las reglas de la sección de agentes del capítulo~\ref{cap:metodologia}. La tabla~\ref{tab:parametros} resume los parámetros.

\begin{table}[htbp]
\centering
\caption{Parámetros del modelo ejecutado (versión 0.4.0).}
\label{tab:parametros}
\small
\begin{tabular}{@{}lr@{}}
\toprule
Parámetro & Valor \\
\midrule
Reservas iniciales $x_0=y_0$ & 10\,000 \\
Precio inicial de la referencia & 1 \\
Horizonte & 20 ticks \\
Tick del choque & 1 \\
Semiancho de la banda $\gamma$ & 0.003 \\
Ticks dentro de la banda para el retorno sostenido & 2 \\
Capacidades de arbitraje $c$ & 0; 0.005; 0.03; 0.10 \\
Inventario de ruido y tendencia (múltiplo de las reservas) & 2 \\
Inventario del arbitrajista (múltiplo) & 5 \\
Probabilidad de actividad del ruido & 0.7 \\
Tamaño de la orden de ruido (fracción de la reserva inicial) & 0.004 \\
Umbral de tendencia (retorno logarítmico) & 0.001 \\
Tamaño de la orden de tendencia & 0.006 \\
Umbral de pánico (retorno logarítmico) & $-0.015$ \\
Tamaño de la venta de pánico & 0.02 \\
Ventas de pánico máximas por ejecución & 1 \\
Disparador del retiro por umbral (error logarítmico) & 0.02 \\
Fracción de reservas que se conserva & 0.55 \\
\bottomrule
\end{tabular}
\end{table}

\section{Contrato de ejecución y reproducibilidad}
\label{sec:contrato}

Cada ejecución del simulador queda descrita por los campos de la tabla~\ref{tab:contrato}.

\begin{table}[htbp]
\centering
\caption{Contrato de ejecución.}
\label{tab:contrato}
\small
\begin{tabularx}{\textwidth}{@{}l X@{}}
\toprule
Campo & Contenido \\
\midrule
\code{run\_id} & Identificador derivado de los parámetros. \\
Commit & Versión del código y estado del árbol de trabajo. \\
Escenario & Condición, capacidad y regla de liquidez. \\
Semilla & Semilla del flujo de ruido. \\
Parámetros & Configuración completa del modelo y del protocolo. \\
Reloj & Tick y evento registrado. \\
Series & Precios, reservas, errores, órdenes. \\
Métricas & Las de la tabla~\ref{tab:metricas}. \\
Comprobaciones & Conservación de activos, positividad y finitud. \\
Entorno & Versión de Python y de la plataforma. \\
Suma de verificación & SHA-256 de configuraciones, código y productos. \\
\bottomrule
\end{tabularx}
\end{table}

\subsection*{Registro de la ejecución Monte Carlo archivada}

\begin{itemize}
  \raggedright
  \item Identificador: \code{\RunID}, versión de análisis 0.5.0, versión del modelo 0.4.0, Python 3.12.3.
  \item Commit de código: \ifdefined\blindreviewmode\code{[omitido para revisión anónima]}\else\code{06778ccfb4db032d8ae143dcd4396d8786ee7814}\fi, con el árbol de trabajo limpio (\code{git\_working\_tree\_dirty} falso).
  \item Recuento de ejecuciones: 3000 de exploración y 13\,500 ejecuciones de simulación reservadas (denominadas históricamente de confirmación).
  \ifdefined\blindreviewmode
  \item Las configuraciones, fuentes, ejecuciones y productos cuentan con sumas de verificación archivadas; sus valores se omiten en esta copia porque permiten localizar los artefactos públicos.
  \else
  \item SHA-256 de la configuración de confirmación archivada:\newline\code{dcf2ef01fdc7504834362f7cb38b93801b2effa861f6d0e3be8ed74796cef16f}; de la configuración del modelo:\newline\code{9c221ada2b8a2150d238c11d320e34fd3697a6a67d193ab4951b48547b09623f}.
  \item SHA-256 del código fuente: \code{confirm.py} \code{9b2c7df09f49bd2f\ldots}, \code{abm.py} \code{f15c666856519b8\ldots}, \code{run.py} \code{1530471be964587e\ldots}, \code{h1c.py} \code{6c26ed25f8796d1c\ldots} (valores completos en el resumen de la ejecución).
  \item SHA-256 de los productos: \code{run\_metrics.csv} \code{2a0fbd17bb3e5eb6\ldots}; \code{inference.csv} \code{d4828a48cd0f8f2a\ldots}; \code{cell\_breakdown.csv} \code{0a87125ca4a09068\ldots}; \code{h3\_by\_condition.csv} \code{28d31f52e5c4be41\ldots};\newline\code{condition\_summary.csv} \code{ee6c02c5ac37544b\ldots};\newline\code{exploration\_fits.csv} \code{1f5f7ffbe00a86c6\ldots}; \code{sample\_size\_guidance.csv} \code{86d4a3a45fedd0a2\ldots}.
  \fi
\end{itemize}

\subsection*{Sensibilidad de ingeniería y repetición completa}

\ifdefined\blindreviewmode
Las fuentes de los ejecutores originales y la configuración archivada no cambiaron. Cinco flujos existentes reprodujeron byte por byte sus artefactos registrados y pasaron las nueve comprobaciones automáticas de implementación y reglas. Los identificadores y hashes que permiten localizar el repositorio se omiten en esta copia.
\else
Las fuentes de los ejecutores originales y la configuración archivada no cambiaron. Cinco flujos existentes reprodujeron byte por byte sus artefactos de resultados registrados: mecánica determinista \code{cex-dex-mvp-f183520b90aa}, barrido \code{cex-dex-sweep-5e33d51c51a3}, selección estocástica \code{cex-dex-abm-a99f9be866b2}, mapa reservado \code{cex-dex-h1c-43e446f04adf} y la ejecución históricamente denominada de confirmación \code{cex-dex-confirm-5b06486d695e}. Pasaron las nueve comprobaciones automáticas de implementación y reglas.
\fi

El análisis suplementario es la sensibilidad de ingeniería (análisis 0.6.0; modelo 0.4.0), 300 semillas nuevas (2000--2299), 40\,500 trayectorias. El protocolo y la configuración documentan tres variantes del flujo de órdenes y las 15 celdas evaluadas. La unidad independiente es la semilla; los intervalos de recuperación por celda son Wilson y los intervalos pareados por profundidad son aproximaciones normales sin ajuste. La reproducción en una segunda carpeta generó artefactos numéricos y de configuración idénticos.

\ifdefined\blindreviewmode
Los comandos y rutas específicos para regenerar la ejecución se omiten en esta copia de evaluación anónima.
\else
Los productos de la ejecución no se versionan en el repositorio. Para regenerarlos y exportar los datos de las figuras:

\begin{verbatim}
python3 experiments/cex_dex_mvp/confirm.py
python3 experiments/cex_dex_mvp/export_thesis_data.py \
  --run-dir experiments/cex_dex_mvp/output/confirm/<run_id>
python3 experiments/cex_dex_mvp/check_invariants.py
\end{verbatim}
\fi

La repetición de la ejecución produjo el mismo identificador y archivos de resultados idénticos byte a byte, salvo los campos \code{generated\_at\_utc} y \code{git\_working\_tree\_dirty} del resumen.

\subsection*{Diagnóstico de razones de capacidad 0.8.0}

La comprobación dirigida repitió 1\,800 trayectorias E3 poco profundas con las semillas archivadas 2000--2299: 900 con límite de arbitraje fijo en las reservas iniciales y pánico activo, y 900 con el mismo límite y pánico desactivado. En la variante con pánico, las tres celdas de choque produjeron 123, 101 y 53 trayectorias con salida, cada una con un solo episodio; los 277 episodios comenzaron en el tick del pánico. La variante sin pánico produjo cero salidas posteriores en las tres celdas (0/900 por celda), y todas las trayectorias se recuperaron con seguimiento posterior a la confirmación. La comprobación es posterior y reutiliza semillas; respalda el diagnóstico configurado $R$ como hipótesis para un factorial preregistrado, no como validación independiente ni umbral universal.

\ifdefined\blindreviewmode
El identificador de ejecución y las sumas de verificación de fuentes y productos se omiten en esta copia anónima.
\else
Las referencias de reproducibilidad son:
\begin{itemize}
  \item Ejecución: \code{cex-dex-capacity-ratio-e269f980fd26}.
  \item Protocolo: \code{experiments/cex\_dex\_mvp/H2\_CAPACITY\_RATIO\_PROTOCOL\_0.8.0.md}.
  \item Script: \code{analyze\_capacity\_ratio\_check.py}.
  \item Manifiesto: \code{figures/data/fixed\_capacity\_ratio\_manifest.json}.
\end{itemize}
\fi

\section{Plan de verificación}
\label{sec:verificacion}

La tabla~\ref{tab:verificacion} lista las nueve comprobaciones automáticas de \code{check\_invariants.py}. Todas pasaron en la versión usada para los resultados archivados. La lista combina comprobaciones mecánicas de implementación y comprobaciones de que ciertos comportamientos esperados se siguen de las reglas configuradas; aprobarlas no constituye evidencia independiente de efectos de mercado.

\begin{table}[htbp]
\centering
\caption{Verificaciones del simulador y resultado.}
\label{tab:verificacion}
\small
\begin{tabularx}{\textwidth}{@{}l X l@{}}
\toprule
Id. & Comprobación & Resultado \\
\midrule
V01 & La venta de $X$ sin comisión conserva $k$ y mantiene reservas positivas. & pasa \\
V02 & La compra de $X$ sin comisión conserva $k$ y mantiene reservas positivas. & pasa \\
V03 & El retiro proporcional conserva el precio marginal. & pasa \\
V04 & Los estados configurados de E0 a E3 tienen reservas positivas y finitas. & pasa \\
V05 & La regla de arbitraje configurada acerca el precio de la piscina a la referencia en ambas direcciones del choque. & pasa \\
V06 & El umbral fijo de retiro se activa con los choques ejecutados y aplica el retiro proporcional configurado. & pasa \\
V07 & La comparación estática y de retiro por umbral comparte la trayectoria exógena y los ticks. & pasa \\
V08 & El impacto de una venta fija crece con cada reducción de liquidez. & pasa \\
V09 & Las condiciones pareadas comparten las órdenes de ruido y conservan los activos (24 trayectorias, 504 estados). & pasa \\
\bottomrule
\end{tabularx}
\end{table}

Estas comprobaciones verifican la implementación. La validación frente a mercados reales no se realizó. No hay una comprobación automática de que el simulador reproduzca hechos estilizados, ni se verificó el comportamiento con comisión, porque no está implementada.

\section{Protocolo de literatura y auditoría de referencias}
\label{sec:protocolo-literatura}

La fuente BibTeX activa contiene 87 registros. La bibliografía compilada de esta revisión contiene 56 elementos citados: 55 trabajos académicos y un registro del repositorio. La verificación de metadatos frente a editoriales (Crossref) o repositorios de preprints cubrió título, autores, año, publicación, volumen y páginas de los trabajos citados. Los trabajos académicos se priorizaron por su relevancia: los centrales sustentan afirmaciones principales (mecánica del AMM, arbitraje, modelos basados en agentes y validación), los recomendados amplían el contexto y los de reserva apoyan una afirmación específica. Los siguientes casos se declaran para hacer explícito su estado:

\begin{itemize}
  \item Campbell et al., Bergault et al., Dean et al. y Cohen et al. se citan como preprints o documentos de trabajo, sin revisión por pares localizada.
  \item En Dean et al., el registro de Crossref indica a Dean como primera autora, y otra base de datos bibliográfica indica otro nombre; se usa Crossref y se recomienda confirmarlo en la página del editor del documento de trabajo.
  \item Windrum et al. (2007) se localizó solo por URL y no tiene DOI verificado en el registro consultado.
  \item Schär (2021): las páginas se infirieron del registro y no se confirmaron en el sitio de la revista. Angeris et al. (2020): no se confirmaron volumen y número.
\end{itemize}

La lista detallada, con la evidencia por entrada, se conserva \ifdefined\blindreviewmode en los materiales del estudio.\else en el repositorio (\code{thesis/literature-audit/}).\fi\ No se encontró un ABM previo de un AMM con retiro de liquidez dependiente del error de seguimiento; esta afirmación se limita a las bases consultadas.

\section{Protocolo de sensibilidad}
\label{sec:sensibilidad}

El trabajo de sensibilidad posterior debe variar los factores candidatos de la tabla~\ref{tab:factores}. Los cambios de un factor a la vez pueden servir de cribado; los diseños factoriales o de cribado \cite{morris1991factorial,saltelli2010variance} son preferibles si el objetivo es estudiar interacciones entre varios factores. Antes de ejecutar un nuevo experimento, se deben especificar el estimando principal, el contraste, la unidad de remuestreo por semilla, el objetivo de precisión práctica y la regla de parada; luego se registra la justificación del tamaño de muestra con el protocolo y los resultados. El archivo histórico \code{figures/data/sample\_size\_guidance.csv} pertenece a un diseño anterior y no es una guía general de tamaño de muestra.

\ifdefined\blindreviewmode
El análisis del mecanismo H2 de la sección~\ref{cap:resultados} repite 9\,000 trayectorias con semillas aleatorias existentes. Incluye trayectorias pareadas con y sin pánico y un control de referencia de capacidad con profundidad baja; se cotejaron 7\,200 trayectorias con límite ligado a reservas frente a métricas archivadas. La repetición es exploratoria, no aporta una muestra independiente ni validación externa, y no identifica causalidad dentro del tick. En esta copia anónima se omiten rutas e identificadores de ejecución.
\else
El análisis del mecanismo H2 de la sección~\ref{cap:resultados} se regenera con \code{python3} \url{experiments/cex_dex_mvp/analyze_h2_revision.py}. La repetición \code{\HtwoReviewRunID} usa 9\,000 trayectorias con semillas existentes e incluye corridas pareadas con y sin pánico y un control de referencia de capacidad con profundidad baja. El script escribe CSV por celda y un manifiesto con sumas de verificación de entradas y productos; 7\,200 trayectorias con límite ligado a reservas se cotejan con métricas archivadas. El análisis es exploratorio, no aporta una muestra independiente ni validación externa, y los episodios al final del tick no identifican causalidad intratick.
\fi

\end{document}
